\documentclass[a4paper]{article}
\usepackage{../mahnigCEC}
\usepackage{programs}
\usepackage{epsfig}
\usepackage{url}
%\input{definition.tex}

\title{MATEDA: A suite of EDA programs in Matlab}
\author{
{\bf{Research Report  EHU-KZAA-IK-2/09}} \\
{\bf{Department of Computer Science and Artificial Intelligence}}\\
{\bf{ University of the Basque Country}}\\ \\
Roberto~Santana$^\ast$, Carlos Echegoyen$^\dagger$, Alexander Mendiburu$^\dagger$,\\
 Concha Bielza$^\ddagger$, Jose A. Lozano$^\dagger$, Pedro~Larra\~naga$^\ddagger$,\\
 Rub\'en Arma\~nanzas$^\dagger$ and Siddartha Shakya$^\star$\\\\
$^\ast$Universidad Polit\'ecnica de Madrid \\
$^\dagger$Intelligent Systems Group \\
  Department of Computer Science and  Artificial Intelligence \\
    University of the  Basque Country \\
$^\ddagger$ Departamento de Inteligencia Artificial,\\
 Universidad Polit\'ecnica de Madrid. Campus de Montegacedo sn. \\
 28660. Boadilla del Monte, Madrid, Spain.\\
$^\star$Intelligent Systems Research Centre, BT Innovate, Ipswich, UK. \\
 {\bf{email}}:roberto.santana@upm.es}% <-this % stops a space

%\pagestyle{empty}

\date{}
\begin{document}
\maketitle

\begin{abstract}                   
  This paper describes MATEDA-2.0, a suite of programs in Matlab for estimation of distribution algorithms.
   The package allows the optimization of single and multi-objective
    problems with estimation of distribution algorithms (EDAs)
  based on undirected graphical models and Bayesian networks. The
  implementation is conceived for allowing the incorporation by the user
  of  different combinations of selection, learning, sampling, and local
  search procedures. Other included methods allow the
  analysis of the structures learned by the probabilistic models,
  the visualization of particular features of these structures and the use of the probabilistic models as fitness modeling tools.
\end{abstract}

\tableofcontents{}

\newpage

\section{Introduction}

 EDAs \cite{Larranhaga_and_Lozano:2002,Lozano_et_al:2005,Muhlenbein_and_Paas:1996,Pelikan_et_al:2006} are evolutionary algorithms based
 on estimation and sampling from probabilistic
 models and able to overcome some of the drawbacks exhibited by
 traditional genetic algorithms (GAs)
  \cite{Goldberg:1989,Holland:1975}. Additionally, the probabilistic models used by EDAs
   can represent
   a priori information about the problem
 structure, allowing a more efficient search of optimal solutions. Learning algorithms can also be used to reveal
 previously unknown information about the structure of black box optimization problems.

  Several EDAs have been proposed for discrete, continuous and mixed
  problems. These algorithms mainly differ in the class of
  probabilistic models employed and the learning and sampling
  methods they use. However, it has been acknowledged that the selection and replacement
  strategies used can also determine important differences in the EDAs behavior.
   Sometimes, it is difficult to decide which is the best EDA choice for a given
   problem and the user would like to compare at least a short list of
   combinations of EDA operators in terms of their efficacy and
   efficiency. Other times, it would be convenient to modify only one
   component of the EDA, keeping the rest of the components
   unchanged. MATEDA-2.0 is conceived for these situations. The idea
   is that the user can easily evaluate a number of variants of EDAs
   before taking a decision about the final implementation.

    Another important use of EDAs is to reveal previously unknown
   information about the search space that has been captured during
   the execution of the algorithm. This information can be encoded
   in the probabilistic models learned or in the points visited
   during the search. Similarly, the probabilistic models can be employed as models of
   functions.  It is useful the design of methods that permit to evaluate the quality
   of the probabilistic models as function predictors.


  MATEDA-2.0  includes three main modules integrated by programs that intend
to fulfill the following objectives:


\begin{itemize}
  \item Optimization module: Implementation of different EDAs for
  single and multi-objective problems.

  \item Data analysis and visualization module: For detecting, extracting and
  visualizing characteristic features in the structures of the models learned during the
  evolution.

 \item Function approximation module: Creation and validation of models of the functions
  based on  the probabilistic models learned by EDAs.

%  \item Analysis of the experiments module: For producing summaries
%  about the performance of the EDAs, including results of
%  statistical tests.
\end{itemize}




\section{Installing MATEDA-2.0}

  MATEDA-2.0 employs the Matlab Bayes Net (BNT) toolbox
\cite{Murphy:2001} and the BNT structure learning package
\cite{Leray_and_Francois:2004}. These programs, which are freely
available from the authors website\footnote{They can be respectively
downloaded from \url{http://people.cs.ubc.ca/~murphyk/Software/BNT/bnt.html}
 and \url{http://banquiseasi.insa-rouen.fr/projects/bnt-slp/}},
should be installed previously to the MATEDA-2.0 installation. Some
of the MATEDA-2.0 routines also employs the MATLAB statistical
toolbox and the affinity propagation  clustering
algorithm\cite{Frey_and_Dueck:2007}\footnote{The Matlab
implementation of affinity
  propagation is available from \url{http://www.psi.toronto.edu/affinitypropagation/}}.

In order to install the software, follow these steps:

\begin{enumerate}
  \item Unpack the file IntEDA.tar.gz and copy the files to a directory
  named MATEDA.
  \item Edit file \texttt{InitEnvironment.m} updating the paths
   'path\_MATEDA', \\ 'path\_FullBNT' and  'path\_BNT\_SLP'.
  \item Set the current Matlab directory to the MATEDA directory.
  \item Execute program \texttt{InitEnvironments.m}. Several warnings  but no error should appear.
\end{enumerate}

The folder \texttt{ScriptMateda.m} contains several  examples of EDAs implementation. If the reader is familiarized with the Matlab environment and EDAs,
these examples should be sufficient for a basic EDA implementation. Otherwise, the following sections
provide a detailed explanation of the  MATEDA-2.0 components.

\section{Defining and executing an EDA}

 The main component of MATEDA-2.0 is the definition of an EDA in
 terms of its parameters and components.  The pseudocode of this  EDA is described in
 Algorithm~\ref{alg:EDA}.



\subsection{General description of the implementation}

 The pseudocode of the general EDA is described in
 Algorithm~\ref{alg:EDA}. Each of the main methods that can be implemented by the user using  MATEDA-2.0 are
 emphasized in Algorithm~\ref{alg:EDA}.


\begin{BAlgo}{Estimation of distribution algorithm}
 \label{alg:EDA}
 \item Set $t\Leftarrow 0$.
  \item \Do
  \item \T{If $t=0$.}
   \item \TT{Generate  an initial population $D_0$ using a \emph{seeding method}.}
   \item \TT {If required, apply a \emph{repairing method} to $D_0$.}
   \item \TT {Evaluate (all the objectives of) population $D_0$ using an \emph{evaluation method}.}
   \item \TT {If required, apply a \emph{local optimization method} to $D_0$.}
   \item \T{Else.}
   \item \TT{Sample a $D_{Sampled}$ population from the model using a \emph{sampling method}.}
   \item \TT {If required, apply a \emph{repairing method} to $D_{Sampled}$.}
   \item \TT {Evaluate (all the objectives of) population $D_{Sampled}$ using an \emph{evaluation method}.}
   \item \TT {If required, apply a \emph{local optimization method} to $D_{Sampled}$.}
   \item \TT{Create a $D_t$ population from populations $D_{t-1}$, $D_{Sampled}$, and  $D_t^S$ using a \emph{replacement method}}
   \item \T {Select a set $D_t^S$ of points according to a \emph{selection method}.}
   \item \T {Compute a probabilistic model of $D_t^S$ using a \emph{learning method}.}
   \item \T {$t \Leftarrow t+1$}
   \item  \Until{The evaluation of the \emph{termination criteria method} is true.}
\end{BAlgo}


 \subsection{Implementation of a general EDA}

 The general EDA program
 \texttt{RunEDA.m} is called as:

  \texttt{[AllStat,Cache] = RunEDA(PopSize,n,F,Card,cache,edaparams);}

 where the input and output parameters have the following meaning:

 \subsubsection{Input parameters}

 \begin{itemize}
   \item $PopSize$: Population size.
   \item $n$: Number of variables.
   \item $F$: Name of the Matlab file that implements the (possibly
   multiobjective) function.
   \item $Card$: Cardinalities of the variables for the discrete
   problems or range of each variable for the continuous problem.
   \item $cache$: A vector specifying which components of the
   algorithm will be stored. $cache(i) = 1$ determines whether the i-th component of
   EDA ($i=1,2,\dots,5$)  will be saved in each generation and $cache(i) = 0$ otherwise. The five components considered are the following:
\begin{enumerate}
  \item Entire population.
  \item Selected population.
  \item Probabilistic model.
  \item Fitness values of the entire population.
  \item Fitness values of the selected population.
\end{enumerate}
   \item $edaparams$: An array of cells specifying all the components
   and parameters used by the EDA. The i-th row of \texttt{edaparams} has the form:




$\{type\_of\_method, name\_of\_implementation, implementation\_parameters\}$.

$type\_of\_method$ defines an EDA component. It is a string
that can take one of  the following values:

\begin{itemize}
\item 'seeding\_pop\_method'
\item 'sampling\_method',
\item 'repairing\_method'
\item 'local\_opt\_method'
\item 'replacement\_method'
\item 'selection\_method'
\item 'learning\_method'
\item 'statistics\_method'
\item 'verbose\_method'
\item 'stop\_cond\_method'
\end{itemize}

$name\_of\_implementation$ is the name of a Matlab program where the
EDA component has been implemented. It can be added by the user or
be one of the methods included in MATEDA-2.0.
$implementation\_parameters$ is a cell array containing the
parameters used by the program \texttt{name\_of\_implementation}
which are passed to it during the execution of \texttt{RunEDA}.
\end{itemize}


\subsubsection{Output parameters}


\begin{itemize}
   \item $AllStat$: For each generation $k$ the cell array  $AllStat\{k,:\}$ contains the following information:
 \begin{itemize}
     \item  $AllStat\{k,1\}$:  Matrix of $5$ rows and $number\_objectives$ columns. Each row shows information about
                      maximum, mean, median, minimum, and variance values of the corresponding objective in the current population
     \item  $AllStat\{k,2\}$:  Stores the best individual.
     \item  $AllStat\{k,3\}$:  Number of different individuals.
     \item  $AllStat\{k,4\}$:   Matrix of $5$ rows and $n$  columns. Each row shows information about
                       maximum, mean, median, minimum, and variance values of
                       the corresponding variable in the current
                       population.
\item  $AllStat\{k,5\}$: Number of function evaluations until
generation $k$.
\item   $AllStat\{k,6\}$: Matrix with the time, in seconds, spent at the
main  EDA steps, each of the $8$ columns stores the times elapsed in
the following steps: sampling, repairing, evaluation, local
optimization, replacement, selection, learning and  total (which
represents the time spent by the previous $7$ and other EDA
operations).
\end{itemize}
   \item If $cache(i)=1$,  for each generation $k$,  $Cache\{i,k\}$ will store
    the corresponding component of the EDA. The order is the one  presented in the explanation of the input parameter $cache$.
\end{itemize}


\begin{example}[Implementation of a continuous Gaussian UMDA]

\begin{programf}
    PopSize = 500;
    n = 30;
    F = 'sum';
    Card(1,:) = zeros(1,n);
    Card(2,:) = 5*ones(1,n);
    cache  = [0,0,0,0,0];
    edaparams\{1\} =  \{'learning\_method','LearnGaussianUnivModel',\{\}\};
    edaparams\{2\} =  \{'sampling\_method','SampleGaussianUnivModel',\{PopSize,1\}\};
    edaparams\{3\} =  \{'replacement\_method','elitism',\{1,'fitness\_ordering'\}\};
    edaparams\{4\} =  \{'selection\_method','prop\_selection',\{\}\};
    [AllStat,Cache] = RunEDA(PopSize,n,F,Card,cache,edaparams);    \end{programf}
\label{ex:EDAPARAMS} \end{example}

The code shown above is an implementation of a continuous Gaussian univariate
marginal distribution algorithm (G-UMDA)
\cite{Larranhaga_et_al:2000,Larranhaga_et_al:1999} using
proportional selection and elitism of one individual. Initially, the
parameters of the EDA are defined as: population size
($PopSize = 500$), number of variables ($n=30$) and fitness function ($F = 'sum'$), which corresponds to the sum of the
variables values, i.e. $f(x) = \sum_{i=1}^n x_i$. The problem has a
continuous representation and therefore the range of values are defined
by instantiating the two rows matrix $Card$, in such a way that
$\forall i, \, 0 \leq x_i \leq 5$. $cache  = [0,0,0,0,0]$ indicates
that the information about the EDA evolution will not be stored.

$edaparams\{1\}$ defines the learning method used which is
implemented in function \texttt{LearnGaussianUnivModel.m}. This
method learns the mean and variance for each of the variables and
outputs them as the model of the selected population. It does not
receive any other parameter apart from the selected population size.
$edaparams\{2\}$ defines the sampling method used which is
implemented in function \texttt{SampleGaussianUnivModel.m}. It
receives the probabilistic model and samples a new population of
$PopSize$ individuals using the mean and variance description of the
data. SampleGaussianUnivModel.m receives two parameters which are
passed in the array $\{PopSize,1\}$.


$edaparams\{3\}$ defines the replacement method which is implemented
in function \texttt{elitism.m}. To determine the best (elite)
solutions, program \texttt{elitism.m} requires the specification of
the number of elitist solutions and a criterion to order the
solutions. These are passed as parameters in the array
$\{1,fitness\_ordering\}$ where \texttt{fitness\_ordering.m}
is a function that orders the individuals according to their fitness
values. The individual with highest fitness value is the first in the ordering. $edaparams\{4\}$ defines the selection method which is
implemented in function \texttt{prop\_selection.m}.


\section{Probabilistic models used in MATEDA-2.0}
\label{sec:PROBMODEL}

The learning and sampling algorithms used by EDAs depend on the
class of probabilistic models. Usually, instances of the same class
of probabilistic models (e.g Bayesian network, Markov chain, etc.)
are learned in each generation. However, we can think of cases where
a different class of probabilistic model could be learned in each
generation \cite{Santana_et_al:2008c}. Switching the class of models
according to the characteristics of the data to be modeled is a
natural way to introduce adaptation in EDAs. MATEDA-2.0 allows the
user to learn different classes of models in each generation. The
only requirement is that the model learned at generation $t$ be
compatible with the sampling algorithm used at the same generation.

In this section, we present the main types of probability models
implemented in MATEDA-2.0.


\subsection{Factorized distributions}

Some probability distributions can be expressed as a product of
marginal probability distributions, each of which is called a factor.
Factorized distributions, or factorizations, are an effective way to
obtain a condensed representation of otherwise very difficult to
store probability distributions. In the following analysis we focus
on factorizations of discrete distributions.

We can identify two components of the factorization.

\begin{enumerate}
  \item The structure of the factorization, which contains
  information about which variables belong to each of the factors
  and the relationships with other factors.
  \item The parameters of the factorizations, which are the
  parameters of each of the factors. In the case of discrete
  factorizations these are usually the probability values of each
  factor configuration.
\end{enumerate}


MATEDA-2.0 represents a factorization using two components:

\begin{enumerate}
  \item $Cliques$, which represent the variables of each factor, specifying whether they
  are also included in previous factors or have not appeared before.
  \item $Tables$, which contain a probability table for each of the
  factors.
\end{enumerate}

Each row of $Cliques$ is a clique. The first column is the number of
overlapping variables with respect to  previous cliques in
$Cliques$. The second column is the number of new variables. Then,
overlapping variables are listed, and finally new variables are
listed. $Tables\{i\}$ stores the marginal tables for clique $i$.

Example~\ref{ex:MODEL} shows different undirected graphs
(Figure~\ref{fig:FACTGRAPHS}), their associate factorizations and
the corresponding MATEDA-2.0 representation (Cliques are
respectively shown in Equations~\ref{eq:CLIQ_a},~\ref{eq:CLIQ_b}
and~\ref{eq:CLIQ_c}).

As illustrated by Example~\ref{ex:MODEL}, the data format used by
$Cliques$ serves to store the wide class of ordered factorizations
\cite{Santana:2005}, which includes marginal product factorizations
\cite{Harik:1999,Muhlenbein_and_Paas:1996,Santana_et_al:2008},
Markov chain factorizations \cite{Santana_et_al:2004},
factorizations constructed from junction trees
\cite{Muhlenbein_et_al:1999,Ochoa_et_al:1999} and junction graphs
\cite{Santana_et_al:2008a}.


\begin{figure*}
\begin{center}
\begin{pspicture}(0,-3.0)(11.4,3.0)
\psset{xunit=.6cm,yunit=.6cm}

\rput(2.5,-5.0){a)} \rput(9.5,-5.0){b)} \rput(16.5,-5.0){c)}

\rput(0,-4){\circlenode{P7}{6}} \rput(5,-4){\circlenode{P8}{7}}
\rput(0,-2){\circlenode{P5}{4}} \rput(5,-2){\circlenode{P6}{5}}
\rput(0,0){\circlenode{P2}{2}} \rput(5,0){\circlenode{P3}{3}}
\rput(2.5,2){\circlenode{P1}{1}}


\rput(7,-4){\circlenode{P7}{6}} \rput(12,-4){\circlenode{P8}{7}}
\rput(7,-2){\circlenode{P5}{4}} \rput(12,-2){\circlenode{P6}{5}}
\rput(7,0){\circlenode{P2}{2}} \rput(12,0){\circlenode{P3}{3}}
\rput(9.5,2){\circlenode{P1}{1}}



\ncline{-}{P5}{P6} \ncline{-}{P1}{P2} \ncline{-}{P1}{P3}
\ncline{-}{P5}{P2} \ncline{-}{P6}{P3} \ncline{-}{P7}{P8}
\ncline{-}{P2}{P3} \ncline{-}{P2}{P6} \ncline{-}{P5}{P3}
\ncline{-}{P3}{P4}

\rput(14,-4){\circlenode{P7}{6}} \rput(19,-4){\circlenode{P8}{7}}
\rput(14,-2){\circlenode{P5}{4}} \rput(19,-2){\circlenode{P6}{5}}
\rput(14,0){\circlenode{P2}{2}} \rput(19,0){\circlenode{P3}{3}}
\rput(16.5,2){\circlenode{P1}{1}}

\ncline{-}{P5}{P6} \ncline{-}{P1}{P2} \ncline{-}{P1}{P3}
\ncline{-}{P5}{P2} \ncline{-}{P6}{P3} \ncline{-}{P2}{P3}
\ncline{-}{P2}{P6} \ncline{-}{P4}{P6} \ncline{-}{P6}{P8}
\ncline{-}{P5}{P7}

%\nccurve[angleB=180]{A}{B}
\end{pspicture}
\caption{Graph representing factorizations of different complexity.
a) Univariate factorization.  b) Factorization with maximum clique
of size $4$. c) Factorization with maximum clique of size $3$.}
\label{fig:FACTGRAPHS}
\end{center}
\end{figure*}



\begin{example}[Representation of different factorizations in MATEDA-2.0]
\begin{align}
p^a({\bf{x}}) &=  \prod_{i=1}^7 p(x_i) \label{eq:fact-a}\\
p^b({\bf{x}}) &= \frac{p(x_1,x_2,x_3) p(x_2,x_3,x_4,x_5)
  p(x_6,x_7)}{p(x_2,x_3)} \label{eq:fact-b}\\
p^c({\bf{x}}) &= \frac{p(x_1,x_2,x_3) p(x_2,x_3,x_5)
  p(x_2,x_4,x_5)p(x_4,x_6) p(x_5,x_7) }{p(x_2,x_3)p(x_2,x_5)p(x_4)p(x_5)} \label{eq:fact-c}
\end{align}



\begin{equation}
 Cliques = \left (\begin{array}{ccc}
           0  & 1 & 1 \\
           0  & 1 & 2 \\
           0  & 1 & 3 \\
           0  & 1 & 4 \\
           0  & 1 & 5 \\
           0  & 1 & 6 \\
           0  & 1 & 7
          \end{array}  \right ) \label{eq:CLIQ_a}
\end{equation}


\begin{equation}
             Cliques =  \left (\begin{array}{cccccc}
             0 & 3 & 1 & 2 & 3 & 0 \\
             2 & 2 & 2 & 3 & 4 & 5 \\
             0 & 2 & 6 & 7 & 0 & 0
           \end{array}  \right )\label{eq:CLIQ_b}
 \end{equation}



 \begin{equation}
Cliques =  \left (\begin{array}{ccccc}
             0 & 3 & 1 & 2 & 3 \\
             2 & 1 & 2 & 3 & 5 \\
             2 & 1 & 2 & 5 & 4 \\
             1 & 1 & 4 & 6 & 0 \\
             1 & 1 & 5 & 7 & 0
           \end{array} \right )\label{eq:CLIQ_c}
 \end{equation}


 \label{ex:MODEL} \end{example}


 Figures~\ref{fig:FACTGRAPHS} a) b) c) show the graphs that correspond to the
 structure of the factorizations represented by
 Equations~\eqref{eq:fact-a},~\eqref{eq:fact-b} and
 ~\eqref{eq:fact-c} respectively. The MATEDA-2.0 representation of these
 factorizations are respectively shown in~\eqref{eq:CLIQ_a},~\eqref{eq:CLIQ_b} and~\eqref{eq:CLIQ_c}.





\subsection{Bayesian and Gaussian networks}

\begin{figure*}
\begin{center}
\begin{pspicture}(0,0)(11,5)

\rput(2,1){\circlenode{P6}{6}} \rput(0,5){\circlenode{P2}{2}}
\rput(1,3){\circlenode{P3}{3}} \rput(3,3){\circlenode{P4}{4}}
\rput(4,5){\circlenode{P5}{5}} \rput(2,5){\circlenode{P1}{1}}


\ncline{->}{P2}{P3} \ncline{->}{P3}{P4} \ncline{->}{P3}{P6}
\ncline{->}{P3}{P1} \ncline{<-}{P4}{P1} \ncline{<-}{P4}{P6}
\ncline{<-}{P4}{P5}

\rput(9,1){\circlenode{P6}{6}} \rput(7,5){\circlenode{P2}{2}}
\rput(8,3){\circlenode{P3}{3}} \rput(10,3){\circlenode{P4}{4}}
\rput(11,5){\circlenode{P5}{5}} \rput(9,5){\circlenode{P1}{1}}


\ncline{->}{P2}{P3} \ncline{->}{P3}{P4} \ncline{<-}{P4}{P1}
\ncline{->}{P3}{P6} \ncline{<-}{P4}{P5} \rput(1.5,0){a)}
\rput(9.5,0){b)}
%\nccurve[angleB=180]{A}{B}
\end{pspicture}
\caption{Structures of: a) General Bayesian network; b) Polytree.}
\label{fig:FACT_BAY_POL}
\end{center}
\end{figure*}

\subsubsection{Bayesian networks}

Bayesian networks are graphical models based on directed acyclic
graphs. They have been used for probabilistic
inference in domains such as expert systems
\cite{Dawid:1992,Lauritzen_and_Spiegelhalter:1988} classification
problems \cite{Blanco:2005,Friedman_and_Goldszchmidt:1996}, and
optimization
\cite{Etxeberria_and_Larranhaga:1999,Pelikan_et_al:1999}.

In a Bayesian network, where the discrete variable $X_i$ has $r_i$ possible
values, $x_i^1, \ldots, x_i^{r_i}$, the local distribution
$p(x_i \mid  {\bf{pa}}_i^{j,S}, {\bf{\theta}}_i)$ is an unrestricted
discrete distribution:

\begin{equation}
p(x_i^k \mid  {\bf{pa}}_i^{j,S},{\bf{\theta}}_i) =
{\bf{\theta}}_{x_i^k \mid {\bf{pa}}_i^{j}}  \equiv \theta_{ijk}
\end{equation}
where ${\bf{pa}}_i^{1,S} , \ldots,{\bf{pa}}_i^{q_i,S}$ denotes the
values of ${\bf{Pa}}_i^{S}$, the set of parents of  ${X_i}$ in the
structure $S$. $q_i$ is the number of possible different instances
of the parent variables of $X_i$, hence $q_i = \prod_{X_g \in
  {\bf{Pa}}_i^{S}}  r_g$.

 The local parameters are given by ${\bf{\theta}}_i = ((\theta_{ijk})^{r_i}_{k=1})^{q_i}_{j=1}$. Parameter $\theta_{ijk}$
 is the conditional probability of variable $X_i$ being in its $k$-th
 state given that the set of parents is in its $j$-th configuration.



\subsubsection{Bayesian network toolbox}

MATEDA-2.0 uses the Matlab Bayes Net (BNT) toolbox
\cite{Murphy:2001} and the BNT structure learning package
\cite{Leray_and_Francois:2004} for the implementation of EDAs that
use Bayesian and Gaussian models. We include a very brief
description of the main features of the toolboxes used by MATEDA-2.0.
Further details could be consulted in
\cite{Leray_and_Francois:2004}.


In the BNT toolbox, a directed acyclic graph (dag) is a structure
that serves to represent the structure of a Bayesian network (bnet).
Example~\ref{ex:BNCREAT} shows the different steps in the creation
of a Bayesian network from a given dag. First, the Bayesian network
is initialized using the dag and the cardinality of the variables.
In the second step, the type of each variable is specified (tabular for
discrete values and Gaussian for continuous variables). Finally, the
parameters of the network are learned from the data. The BNT toolbox
includes methods for learning the Bayesian network structure from
the data. They are analyzed in Section~\ref{sec:LEARNMETHODS}.


\begin{example}[Bayesian network learning in BNT]
\begin{programf}
n = 4;
Card = 2*ones(1,n);
dag(1,:) = [0 0 0 1];
dag(2,:) = [0 0 0 0];
dag(3,:) = [0 1 0 1];
dag(4,:) = [0 0 0 1];
init\_bnet = mk\_bnet(dag,Card);
for=1:n
 init\_bnet.CPD{i} = tabular\_CPD(init\_bnet,i);
end,
bnet = learn\_params(init\_bnet,data);\end{programf}
\label{ex:BNCREAT} \end{example}

\subsubsection{Gaussian networks}


In a Gaussian network \cite{Shachter_and_Kenley:1989}, each variable $X_i \in \mbox{\boldmath $X$}$ is continuous and
each local density function is the linear-regression model:

\begin{equation}
f (x{_i} \mid \mbox{\boldmath $pa$}_{i}^{S}, \mbox{\boldmath $\theta$}_i) \equiv  {\mathcal{N}} (x_i; m_i + \sum _{x_{j} \in \mbox{\boldmath $pa$}_{i}} b_{ji} (x_j - m_j), \frac{1}{v_i})
\end{equation}
where ${\mathcal{N}} (x; \mu, \sigma^{2})$ is a univariate normal
distribution with mean $\mu$ and variance $\sigma^{2}$. Given this form, a missing arc from $X_j$ to $X_i$
implies that $b_{ji} = 0$ in the former linear-regression model.
The local parameters are given by  $\mbox{\boldmath $\theta$}_i =
(m_i,  \mbox{\boldmath $b$}_i, v_i)$, where $\mbox{\boldmath
$b$}_i= (b_{1i}, \ldots, b_{i-1i})^t$ is a column vector.

The interpretation of the components of the local parameters is as
follows: $m_i$ is the unconditional mean of $X_i$, $v_i$ is the
conditional variance of $X_i$ given $\mbox{\boldmath $Pa$}_{i}$,
and $b_{ji}$ is a linear coefficient reflecting the strength of
the relationship between $X_j$ and $X_i$.


\subsection{Markov networks} \label{MRF}

Given an undirected graph $G=(V,E)$ we have the following
definitions \cite{Hammersley_and_Clifford:1968}:

\begin{definition}
   The neighborhood $N(X_i)$ of a vertex $ X_i \in {\bf{X}}$ is defined as
  $N(X_i)  = \{ X_j: (X_j,X_i) \in E \}$. The set of edges uniquely
  determines a neighborhood system on the associated graph $G$.
\label{NEIGHBHOOD}
\end{definition}

\begin{definition}
   The boundary of a set of vertices, ${\bf{X}}_S \subseteq {\bf{X}}$,
 is the set of vertices in $
 {\bf{X}} \setminus {\bf{X}}_S$ that neighbors  at least one vertex in ${\bf{X}}_S$. The boundary
 of ${\bf{X}}_S$ is denoted as $bd({\bf{X}}_S)$.
\label{BOUNDARY}
\end{definition}

\begin{definition}
   The closure of a set of vertices, ${\bf{X}}_S \subseteq {\bf{X}}$, is the set of vertices
 $cl({\bf{X}}_S) = {\bf{X}}_S \cup bd({\bf{X}}_S)$.
\label{CLOSURE}
\end{definition}


\begin{definition}
 A probability $p({\bf{x}})$ is called a Markov random field with respect to the neighbor
system on  a graph $G$ if, $\forall   {\bf{x}} \in {\bf{X}}, \;
\forall i \in \{1, \ldots, n\}$,  $p(x_i|{\bf{x}} \setminus
x_i)=p(x_i|bd(x_i))$.
\end{definition}

\begin{definition}
A probability $p({\bf{x}})$ on a graph $G$ is called a Gibbs field
with respect to the neighborhood system on the associated graph $G$
when it can be represented as follows:
\begin{equation}
 p({\bf{x}}) = \frac{1}{Z} e^{-H({\bf{x}})}  \label{GIBBSAMP}
\end{equation}
 where  $H({\bf{x}}) = \sum_{C \in \mathcal{C}} \Phi_{C}({\bf{x}})$
 is called the energy function,  being $\Phi = \{\Phi_{C} \in C \}$  the set of clique potentials, one for each of the  maximal
cliques in $G$. The value of $\Phi_{C}({\bf{x}})$  depends on its
local configuration on the clique $C$. The normalizing constant $Z$
is the corresponding partition function,  $Z =\sum_{{\bf{x}}}
e^{-H({\bf{x}})}$.
\end{definition}

 In MATEDA-2.0, $Cliques$ are used to represent the neighborhood structure
in models based on Markov networks
\cite{Santana:2003c,Shakya_and_McCall:2007,Shakya_and_Santana:2008}.
In this particular case,   the first column of $Cliques(i,:)$
represents the number of neighbors for variable $X_i$.  The second,
is the number of new variables (for Markov networks, only one new
variable $X_i$ appears in each clique.).  Then, the neighbor variables are
listed and finally variable $X_i$ is added.

 The parameters of the Markov network are represented using the
 $Tables$ structure which stores the conditional probabilities of each variable given its
 neighbors.


\subsection{Mixtures of distributions}

 A mixture of distributions \cite{Everitt_and_Hand:1981,McLachlan_and_Basford:1988} is defined to be a distribution of the form:

\begin{equation}
   Q(x)=\sum\limits_{j=1}^m\lambda _j \cdot f_j(x)  \label{eq:Mixture}
\end{equation}

with $\lambda _j > 0$, $j=1,\dots,m$, $\sum_{j=1}^m\lambda _j=1$.


$f_j$ are called component densities or  mixture components, and
$\lambda _j$ are called mixture proportions or mixture
coefficients. $m$ is the number   of components of the mixture. A
mixture of distributions can be viewed as containing an unobserved
choice variable $z$, which takes values $k\in $ $\{1,\dots,m\}$ with
probability $\lambda _j$. In some cases the choice variable $z$ is
known.

 Examples of mixtures distributions include:
 \begin{itemize}
  \item Mixtures of Gaussian distributions \cite{Everitt_and_Hand:1981}.
  \item Bayesian multi-nets \cite{Geiger_and_Heckerman:1996}.
  \item Mixtures of trees \cite{Meila_and_Jordan:2000}.
 \end{itemize}

  In EDAs, mixtures of distributions have been extensively applied
  to solve discrete \cite{Bosman_and_Thierens:2002,Pelikan_and_Goldberg:2000,Penha_et_al:2002,Santana_et_al:2001b,Santana_et_al:2006a}
   and continuous \cite{Bosman_and_Thierens:2002,Cho_and_Zhang:2002,Gallagher:2000,Thierens_and_Bosman:2001a} optimization problems.

 In Mateda-2.0, a mixture of models is represented as an array of
 components and coefficients. Each element of the array stores all
 the relevant information about the corresponding model. A mixture
 can contain as components models of different classes.
 Example~\ref{ex:MIXTURE} shows the way in which a mixture of
 Gaussian distributions can be learned from a cluster of solutions.

\begin{example}[Constructing a mixture of Gaussian distributions]
\begin{programf}
 \% nclusters: number of clusters where solutions in AuxPop are
 grouped
 \% ind: index of the cluster each solution belongs to.
 for i=1:nclusters,
   idx = find(ind==i);
   nmembers = size(idx,1); \% Number
   model{1,i} =  mean(AuxPop(idx,:)); \% Vector of means for each cluster
   model{2,i} =  std(AuxPop(idx,:));  \% Vector of standard deviation for each cluster
   model{3,i} =  nmembers/PopSize;    \% Coefficient for the mixture, proportional to the
 end                  \UnnumLine         \% number of points in the cluster
   \end{programf}
\label{ex:MIXTURE} \end{example}


\subsection{Non probabilistic models}

 Even if MATEDA-2.0 is conceived for EDAs where modeling of the
 solutions is done using probabilistic models, there may be
 situations in which population based algorithms that use other types of modeling
 techniques (e.g. algorithms  that use modeling based on neural networks)
 or do not use any type of modeling (e.g. genetic algorithms \cite{Goldberg:1989}) are to be compared with EDAs.
 Assuming that the main difference of these algorithms with EDAs lies in the variation operation they use
 (i.e. the evaluation and selection steps are essentially identical), these algorithms
 could be implemented in MATEDA-2.0 using two alternatives:

 \begin{itemize}
  \item Split the implementation of the variation operator in two
  parts. In the first part, a \emph{model} is learned. In the
  second,  the model is used to sample solutions.
  \item Define a learning method where nothing is done and implement
  the variation operators in the sampling method.
 \end{itemize}

 Example~\ref{ex:CXMODEL} presents a model of one-point crossover operator of a GA.
 The models correspond to a list of triples, where the first value
 of each triple corresponds to the crossover point for the creation
 of an individual and the other two values are the indices of the
 parents in the selected population. In this simple example, there will be
 generated $10$ new individuals. The first two offsprings will be formed, during the sampling
 step, by crossing parents $1$ and $3$ at point $5$.

\begin{example}[Model of one-point crossover for a GA]
 \begin{equation}
Cliques =  \left (\begin{array}{ccc}
             5 & 3 & 1  \\
             6 & 4 & 2  \\
             2 & 6 & 2  \\
             8 & 5 & 3   \\
             7 & 2 & 5
           \end{array} \right )\label{eq:CX}
 \end{equation}
 \label{ex:CXMODEL} \end{example}


\section{Implementation of the EDA components}

There is a common way in which all EDA components are invoked from the main program  \texttt{RunEDA} using the Matlab function \texttt{eval}. This
is:

 \texttt{[Output] = eval([eda\_method,'(fixed\_params,user\_params)']);}
where $eda\_method$ is the name of the Matlab program that implements the
 method. $fixed\_params$ are a set of parameters determined by the
 \texttt{RunEDA} implementation and $user\_params$ are parameters defined by
 the users and passed to the program using $eda\_params$. In the
 following section we briefly describe the format of the
implementations for each of the EDA components. Afterwards, several
examples that illustrate the use of these components are discussed.

\subsection{Seeding methods}

 Seeding methods are used to initialize the starting population when
 some knowledge about the problem is available or the user wants to
 bias the search to certain areas of the space.

In  \texttt{RunEDA}, the seeding procedure is invoked at the first
population as:

\texttt{Pop=eval([seeding\_pop\_method,'(n,PopSize,Card,seeding\_pop\_params)']);}

 The seeding methods implemented in MATEDA-2.0 are the following (see help
function-name for details on the input parameters used by the methods):

\begin{itemize}

 \item \texttt{seed\_thispop}: Seeds a given population.

 \item  \texttt{RandomInit}:  Samples a population of solutions from the uniform
distribution.

\item  \texttt{Bias\_Init}:   Biased initialization of a population of binary vectors,    where the probability of generating $1$ is $p$ and the probability of generating $0$ is $1-p$.

 \item  \texttt{seeding\_unitation\_constraint}:  Generates a population of binary vectors
                               where all vectors have the same number of
                               ones. It might be used for problems
                               with unitation constraints \cite{Santana_and_Ochoa:1999b,Santana_et_al:2001c}.


\end{itemize}

By default, $seeding\_pop\_method = RandomInit$.


\subsection{Sampling methods}

 Sampling methods serve to generate the new population from the
 probabilistic model learned by the EDA. The type of sampling method
 is therefore dependent of the class of probabilistic method used.
 Traditionally, probabilistic logic sampling (PLS)
 \cite{Henrion:1988} has been the choice in EDAs, but other
 proposals incorporate Gibbs sampling \cite{Gamez_et_al:2007,Santana:2005,Shakya_and_McCall:2007,Shakya_and_Santana:2008}, and methods that find the most probable
 configurations \cite{Mendiburu_et_al:2007a,Santana:2006,Soto:2003}.
 MATEDA-2.0 implements variants of all these sampling methods.


In  \texttt{RunEDA}, the sampling procedure is invoked at every
generation, except the first one where seeding is applied, as:

\texttt{ NewPop
=eval([sampling\_method,'(n,model,Card,SelPop,SelFunVal \\
,sampling\_params)']); }

where $model$ is a cell array  containing
the description of the probabilistic model. The inclusion of the
selected population and its evaluation as parameters of the sampling
method allow the implementation of sampling algorithms that may
start from previously found solutions (e.g. Gibbs sampling).


The following sampling  methods have been implemented (see help
function-name for details on the  input parameters used by the methods):

\begin{itemize}
  \item \texttt{MOAGeneratePopulation}: Samples a Markov network  using Gibbs
  Sampling.
  \item \texttt{SampleFDA}: Samples  from a factorized model using
  PLS.
  \item \texttt{SampleGaussianUnivModel}: Samples  from a univariate Gaussian
  model.
   \item \texttt{SampleGaussianFullModel}: Samples from a full multivariate Gaussian
  model.
  \item \texttt{SampleMixtureofUnivGaussianModels}:  Samples  from a mixture of univariate Gaussian
  models.
    \item \texttt{SampleMixtureofFullGaussianModels}:  Samples  from a mixture of full multivariate Gaussian
  models.
  \item \texttt{SampleBN}: Samples a  population of individuals from a Bayesian or Gaussian network using PLS.
  \item \texttt{SampleMPE-BN}: Samples a population of discrete solutions in which the first individual corresponds
   to the most probable configuration given by the model and the remaining individuals  are sampled using PLS.
\end{itemize}


The \texttt{SampleBN} and \texttt{SampleMPE-BN} programs invoke
methods defined in the BNT toolbox. EDAs that  use multivariate
Gaussian distributions have shown to suffer from stagnation due to a
very fast (at exponential rate) decrease of the variance
\cite{Bosman_and_Grahl:2008}. A partial remedy is to use an
artificial expansion of the variance.  Therefore, in MATEDA-2.0 some
of the EDAs that use Gaussian models include as a feature a variance
scaling parameter. Other adaptive schedules
\cite{Bosman_et_al:2008a,Dong_and_Yao:2007,Grahl_et_al:2006,Posik:2008}
for modifying the variance during sampling could be implemented in
this framework.

\subsection{Repairing methods}

  Repairing methods are commonly used to guarantee the feasibility
  of solutions. They modify (repair) a given individual to guarantee
  that the constraints are satisfied.

In  \texttt{RunEDA} the repairing procedure is invoked at every
generation  as:


  \texttt{[Pop] =
  eval([repairing\_method,'(Pop,Card,repairing\_params)'])}

The following repairing  methods have been implemented (see help
function-name for details on the functions parameters):

\begin{itemize}
 \item \texttt{SetInBounds\_repairing}: For a problem with continuous representation, this program
 changes the values of each out of range variable  to the minimum (respectively maximum) bounds
 if variables are under (respectively over) the variables ranges.
\item \texttt{SetWithinBounds\_repairing}: For a problem with continuous representation, this program
 truncates the values of each out of range variable to a random
 value within the feasible range.
 \item \texttt{HP\_repairing}:   Recursive procedure, originally introduced in \cite{Cotta:2003},
 to repair a grid path to avoid self-intersections, guaranteeing feasibility.
\end{itemize}



\subsection{Local optimization methods}

Local optimization methods are used to improve the solutions sampled
by EDAs. The combination of EDAs and local optimization methods
\cite{Muehlenbein_and_Mahnig:2002,Pelikan_et_al:2002,Santana_et_al:2008f,Zhang_et_al:2005a}
have been shown to produce important improvements in the efficiency
of the optimization process.


In  \texttt{RunEDA} the local optimization method  is invoked at
every generation as:

  \texttt{  [Pop,NumbEvals] =  eval([local\_opt\_method,'(k,Pop,FunVal,\\
   local\_opt\_params)'])}

$FunVal$ is a matrix with fitness values of all the individuals in
$Pop$ for all the objectives. Since the local optimization program
receives the entire population, it is possible to implement
optimization methods that use information about the population to
modify each individual. Otherwise, the local optimization method can
be applied to each individual separately. $NumbEvals$ is the number
of evaluations made during the local optimization step.


\subsection{Replacement methods}

 Replacement methods are used to combine the individuals generated in the current
 population with individuals from previous generations. In EDAs, they can be
 used as an effective way to promote diversity in the population (e.g. restricted tournament
 replacement \cite{Pelikan:2005}). In addition, research on the role of replacement methods in
 EDAs \cite{Lima_et_al:2007} has shown
 that the type of replacement method can influence the accuracy of the
 probabilistic models learned by the algorithms.   Niching methods
 \cite{Mahfoud:1995} may also be implemented at this step.


In  \texttt{RunEDA} the replacement method  is invoked at every
generation except the first one as:

  \texttt{ [Pop,FunVal]  = eval([replacement\_method, '(OldPop,SelPop,NewPop,\\
   OldFunVal,SelFunVal,NewPopFunVal,replacement\_params)']); }

where $OldPop$, $SelPop$ and $NewPop$ are respectively the
population at step $t-1$, the selected population and the population
generated by sampling (and possibly repaired or improved using local
optimization). $OldFunVal$, $SelFunVal$, and $NewPopFunVal$ are
matrices with  the fitness values, for all the objectives, of the
individuals in $OldPop$, $SelPop$ and $NewPop$.

The following replacement methods have been  implemented:

\begin{itemize}
 \item \texttt{elitism}: Adds the $k$ best individuals of the previous population
to the current population.
 \item  \texttt{best\_elitism}: Joins the selected individuals with the sampled individuals to form the next
population. It should be enforced that $SampledPopSize = PopSize -
SelPopSize$.
 \item \texttt{popaggregation}: Joins the current population with the sampled population
and selects the best $PopSize$ individuals as the new population.
 \item none: $Pop = NewPop$.
 \item \texttt{RT\_replacement}: Creates a new population ($NewPop$) applying restricted tournament replacement.
\end{itemize}

All these methods require the definition of an ordering criterion to
determine the best individuals. While this criterion may be
straightforward for the single objective optimization, there are
many more options available for the case of multi-objective
optimization.

In  \texttt{RunEDA}, when an ordering of the individuals is
required, an ordering method is invoked as:

\texttt{[Ind]  =
eval([find\_bestinds\_method,'(Pop,FunVal)'])}

where $find\_bestinds\_method$ is the name of the ordering method,
passed as a parameter of the main EDA component that uses it (e.g.
replacement method), $Pop$ is a population and $FunVal$ the
corresponding evaluations of the individuals for all the objectives.
The method outputs an index $Ind$ with the ordering.

The following ordering methods have been  implemented:

\begin{itemize}
 \item \texttt{fitness\_ordering}: Individuals are ordered according to the ranking of its fitness values (best is the
 maximum).
 For multi-objective functions, individuals are ordered according to the average ranking for all the objectives (i.e. for each objective an ordering
 is done, and they are averaged later).
 \item \texttt{Pareto\_ordering}:  Individuals are ordered according to the front they belong to. Individuals in
 the first front (nondominated solutions) come first. Then individuals that
 are  only  dominated by those in the first front and so on. There is no ordering criteria
 for individuals in the same front. See \texttt{ParetoRank\_ordering} for more refined
 ordering.
 \item \texttt{ParetoRank\_ordering}: Individuals are firstly ordered according to the front they belong to.
 Then, in each front, they are ordered according to the average rank of their fitness functions. The first front (non-dominated solutions) comes first. Then individuals that are only
 dominated by those in the first front and so on.
\end{itemize}



\subsection{Selection methods}

Selection methods determine which the set of individuals that will
be modeled is. Usually, it is expected that these individuals
correspond to a sample of a current promising region of the search
space. Selection methods are essential for EDAs.  Therefore, some
work has been devoted to study the
 role of selection operators in EDAs
\cite{Johnson_and_Shapiro:2002,Lima_et_al:2007,Mahnig_and_Muehlenbein:2001,Santana_et_al:2005b}.
In  particular, to investigate
 their influence in  the relationship
 interactions-dependencies. It has been shown \cite{Lima_et_al:2007} that while some
 selection methods (e.g. truncation selection) are good at keeping the
 original interactions between the problem variables,  other selection methods can be better in terms of the
 number of function evaluations to reach the optimum (e.g. tournament selection).


In  \texttt{RunEDA} the selection method  is invoked at every
generation as:

  \texttt{[SelPop,SelFunVal]  = eval([selection\_method,'(Pop,FunVal,\\
   selection\_params)'])}

where the input and output parameters are as described for previous
methods.

The following selection methods have been  implemented:

\begin{itemize}
 \item \texttt{exp\_selection}: The selection probability of each individual is exponential  of $base$ to its fitness. The selected population is sampled using stochastic
 universal sampling (SUS) \cite{Baker:1987}. Implemented for single objective functions.
 \item \texttt{prop\_selection}: The selection probability of each individual is proportional to its fitness.
  The selected population is sampled using SUS. Implemented for single objective functions.
 \item \texttt{NonDominated\_selection}: The set of non dominated individuals is selected.
 Implemented for multi-objective functions. The number of selected
 individuals may change at each generation.
 \item \texttt{truncation\_selection}: Individuals are ordered according to the given
 ordering (see previous section for implementations of the ordering
 criterion). The best $T*PopSize$ are selected where $T$ is the truncation coefficient.
\end{itemize}



\subsection{Learning methods} \label{sec:LEARNMETHODS}

  In EDAs, the probabilistic models are constructed using learning
  methods that build the models extracting relevant information
  from the data and (possibly) a priori knowledge about the model
  structure.  Regarding the way learning is done
   in the probability model, EDAs can be divided into two classes \cite{Larranhaga:2002a}. One class groups the algorithms that
   do a parametric learning of the probabilities, and the other one comprises
   those  algorithms where  structural learning of the model is
   also done. Parametric and  structural learning are also known as model fitting
   and model selection, respectively.


In  \texttt{RunEDA} the learning method  is invoked at every
generation as:

  \texttt{[model] = eval([learning\_method,'(k,n,Card,SelPop,\\
  SelFunVal,learning\_params)'])}

where $k$ is the number of the current generation  and $model$ has
the structure explained in Section~\ref{sec:PROBMODEL}. The
structure and parameters contained in $model$ have to be consistent
with the sampling algorithm used in the same generation.

To learn  Bayesian and Gaussian networks, MATEDA-2.0 uses a number
of methods implemented in the Matlab Bayes Net (BNT) toolbox
\cite{Murphy:2001} and the BNT structure learning package
\cite{Leray_and_Francois:2004}.
The following learning methods have
been  implemented:



\begin{itemize}
  \item \texttt{LearnBN}: Learns a Bayesian network from data using the
  following approaches:
  \begin{itemize}
    \item  PC algorithm \cite{Spirtes_et_al:1993}  as implemented in \cite{Murphy:2001}.
    \item  Tree learning algorithm using the mutual information
    \cite{Chow_and_Liu:1968} as implemented in  \cite{Leray_and_Francois:2004}.
    \item Greedy score+search  \cite{Cooper_and_Herskovits:1992}  with a randomly generated initial node
    ordering and Bayesian information criterion (BIC) \cite{Schwarz:1978} scores as implemented in
     \cite{Murphy:2001}.
    \item Greedy score+search  with an initial node
    ordering constructed from a learned tree structure \cite{Leray_and_Francois:2004}, and Bayesian and BIC scores \cite{Cooper_and_Herskovits:1992} as implemented in
     \cite{Murphy:2001}.
  \end{itemize}
  \item \texttt{LearnFDA}:  Creates a factorized model from a given fixed undirected structure.
  \item \texttt{LearnMargProdModel}: Learns a marginal product \cite{Santana_et_al:2008} model using affinity propagation \cite{Frey_and_Dueck:2007} on the matrix of mutual information.
  \item \texttt{LearnMOAModel}: The Markov network model learned by the Markov optimization algorithm (MOA) \cite{Shakya_and_Santana:2008} is
  learned.
  The structure of the model can be given or learned from the data.
  \item \texttt{LearnGaussianUnivModel}: Learns a Gaussian univariate marginal product
  model \cite{Larranhaga_et_al:1999}.
   \item \texttt{LearnGaussianFullModel}: Learns a  full Gaussian multivariate model \cite{Larranhaga_et_al:1999}.
  \item \texttt{LearnMixtureofUnivGaussianModels}: Learns a mixture
  of  univariate Gaussian models using clustering.
   \item \texttt{LearnMixtureofFullGaussianModels}: Learns a mixture of
   full  multivariate Gaussian models using clustering.
  \item \texttt{LearnGaussianNetwork}: Learns a Gaussian network from data using the
  following approaches:
  \begin{itemize}
    \item  Tree learning algorithm using the  BIC score as implemented in
    \cite{Leray_and_Francois:2004}.
    \item Greedy score+search  with a randomly generated initial node
    ordering and  BIC score \cite{Cooper_and_Herskovits:1992} as implemented in
     \cite{Murphy:2001}.
    \item Greedy score+search   \cite{Cooper_and_Herskovits:1992} with  an initial node
    ordering constructed from a learned tree structure \cite{Leray_and_Francois:2004}, and BIC score as implemented in
     \cite{Murphy:2001}.
  \end{itemize}
\end{itemize}


The implemented learning methods incorporate parameters that add
flexibility to the EDAs. For instance, EDAs  that use mixtures of
Gaussian models are based on a clustering of the solutions.
Clustering can be done in the space of variables, the space of
objectives, or both together. User defined clustering methods can be
employed (currently available are affinity propagation and k-means)
with  different distance measures. Similarly, it could be also
possible to add adaptation to learning by setting the coefficients
of the mixture's components according to the accuracy of each
component approximation \cite{Santana:2002}.


Notice that the computational cost in terms of time and memory is very high for some of the Bayesian and Gaussian network learning methods making them
appropriate for only small problems. See file \texttt{ScriptMateda.m} for examples of learning algorithms used for problems of many variables.



\subsection{Methods to compute the statistics}

 To evaluate the performance of an EDA it is necessary to compute a
 number of statistics that describe the behavior of the algorithm.
 MATEDA-2.0 allows the user to implement a procedure that processes
 the information and computes a number of relevant statistics that
 could be used by the methods that display information about the
 algorithm during the evolution.

In  \texttt{RunEDA}, the method to compute the statistics  is
invoked at every generation as:

  \texttt{[AllStat] = eval([statistics\_method,'(k,Pop,FunVal,time\_operations,\\
  number\_evaluations,AllStat,statistics\_params)'])}
where $time\_operations$ is a cell array containing the time (in
seconds) spent by each component of the EDA at each generation,
$number\_evaluations$ is a vector with the number of evaluations in
each generation and  $AllStat$ is an array that contains the
statistics at each generation and is updated by the method for
computing the statistics.

The following method for computing the statistics  has been
implemented:

\begin{itemize}
 \item \texttt{simple\_pop\_statistics}:  Computes the following statistics in
 each EDA
 generation and stores them in $AllStat$.
  \begin{itemize}
   \item   Information about maximum, mean, median, minimum, and variance values of
every objective in the current population.
    \item  Best individual (according to an ordering criterion).
    \item  Number of different individuals.
    \item Information about maximum, mean, median, minimum, and variance values of
every variable in the current population.
  \end{itemize}

\end{itemize}


\subsection{Methods to display information during the evolution}

  An important feedback about the behavior of EDAs is received from
  the type of
  information displayed during the evolution of the optimization algorithm. In
  MATEDA-2.0, this method may be implemented by the user.


In  \texttt{RunEDA}, the method to display the EDA information is
invoked at every generation as:

  \texttt{eval([verbose\_method,'(k,AllStat,verbose\_params,auxedaparams)'])}
where $AllStat$ is the cell array containing the information
computed by the statistics method and $auxedaparams$ contains the
description of all the EDA methods.

The following method for displaying the information about the EDA
behavior has been implemented:

\begin{itemize}
 \item  \texttt{simple\_verbose}: Displays the information computed by the \texttt{simple\_pop\_statistics}. In addition, the
  number of evaluations and the time spent by each EDA component are
  displayed. A description of each the  methods used by the EDA is displayed as a preamble to the
execution of the algorithm.
\end{itemize}

The method used to compute and store the statistics and the one used
to display them should be consistent.


\subsection{Methods to evaluate the stopping conditions}

  The stopping criterion used in EDA influences  the efficiency of the
  algorithm. Several stopping criteria can be used, separately or
  combined. In MATEDA-2.0 the stopping criteria can be defined by the
  user by means of a method that determines whether the algorithm
  should or not stop.

  In  \texttt{RunEDA}, the method to evaluate the stopping conditions is
invoked at every generation as:

\texttt{continue\_evolution
= eval([stop\_cond\_method,'(k,Pop,FunVal,stop\_cond\_params)'])}

The stopping condition criteria will use information about the
current generation, the genotypic structure of the population and
the function values.

The following methods for displaying the information about the EDA
behavior have been implemented:

\begin{itemize}
 \item \texttt{maxgen}:  The algorithm stops when a maximum number of generations has been reached.
 \item \texttt{maxgen\_maxval}:  The algorithm stops either when
 a maximum number of generations has been reached or a given  value of the function has been  reached.
\end{itemize}



\section{Functions and testbed problems implemented in MATEDA-2.0}

To validate the behavior of the EDAs, MATEDA-2.0 includes the
implementation of a number of optimization functions. In this
section we present some of the functions that have been implemented
and a number of methods used to generate instances of these
functions. The use of EDAs to address more realistic problems are
discussed in the next section.

\subsection{Single objective functions}


\subsubsection{Additively decomposable functions}


An additively decomposable function (ADF) is represented as follows:
\begin{equation}
  f({\bf{x}}) = \sum_{i=1}^m f_i({\bf{x}}_{s_i})
\end{equation}
where ${\bf{x}}_{s_i}$ are subvectors  of ${\bf{x}}$ called the
definition sets of the function, and  $f_i$ are
 its defined subfunctions, or just the subfunctions, of
$f({\bf{x}})$.

 ADFs are a clear example of functions where problem structure is known.  In an
ADF, structurally related variables are usually considered as those
that belong to the same definition set. ADFs have been extensively
employed to evaluate EDAs. In particular, deceptive ADFs have been used to illustrate the capacity of
EDAs to deal with the linkage problem affecting simple genetic
algorithms.

  In MATEDA-2.0, an ADF can be implemented as the sum
  of its subfunctions for each definition set.
  The following ADFs have been implemented as examples:

\begin{itemize}
  \item Goldberg deceptive function \cite{Goldberg:1987}.
  \item $Trap$ function of order $k$ \cite{Ackley:1987}.
\end{itemize}

 Single objective functions can also be implemented as particular
 case of multi-objective functions discussed in the next section.


\subsection{Multi-objective functions}

\subsubsection{Multi-objective additive decomposable functions}

 MATEDA-2.0 includes a number of methods that allow the
definition of multi-objective decomposable functions. This class of
functions have been proposed \cite{Khan:2003,Khan_et_al:2002,Pelikan_et_al:2005} to investigate the behavior
 of EDAs for multi-objective problems. In some cases
\cite{Khan:2003,Khan_et_al:2002,Pelikan_et_al:2005}, the rationale
behind the use of these functions has been to obtain objectives that
compete in all or most partitions of an appropriate problem
decomposition.

In MATEDA-2.0, general multi-objective decomposable functions can be
defined by the user by means of the global variables
$FunctionStructure$ and $FunctionTables$. The first global variable
defines which problem variables are involved in the evaluation of
each of the objectives. $FunctionTables$ defines the values of an
objective for each configuration of the variables it depends on.

The main difference between this type of multi-objective
decomposable functions and ordinary ADFs is that each subfunction
will map to a different objective. As a result, we can have
different objectives defined on the same subset of variables.

The following methods have been implemented:

 \begin{itemize}
   \item EvaluateGeneralFunction: Evaluates a vector on a multimodal function whose structure and
          values are respectively defined as global variables $FunctionStructure$ and $FunctionTables$.
   \item PartialEvaluateGeneralFunction:  Evaluates a vector on a multimodal function whose structure and
          values are respectively defined as global variables but only a subset of the objectives are evaluated.
   \item SumEvaluateGeneralFunction: Evaluates a vector on a multimodal function whose structure and
          values are respectively defined as global variables. The Sum of
          the objectives is given as output.
   \item SumPartialEvaluateGeneralFunction: Evaluates a vector on a multimodal function whose structure and
          values are respectively defined as global variables
          but only a subset of the objectives are evaluated and the sum of
          these objectives is given as output.
 \end{itemize}



One multi-objective decomposable function can be transformed to a
single objective function by adding the values of all the
objectives.  This  transformation, that could be seen as the
opposite direction of the so-called \emph{multiobjectivization}
technique \cite{Knowles_et_al:2001} can be done in MATEDA-2.0 using
\texttt{SumEvaluateGeneralFunction} method. Additionally, one may be
interested in evaluating just a subset of objectives. This can
be accomplished by specifying a global variable $SelectedObjectives$
and applying the method $PartialEvaluateGeneralFunction$.

Example~\ref{ex:MULTOBJT} shows the generation and optimization of a
multi-objective decomposable function using MATEDA-2.0.

\begin{example}[Optimization of multi-objective decomposable functions]
 \begin{programf}
 PopSize = 1000;
 n = 50;
 cache = [1,1,1,1,1];
 Card = 2*ones(1,n);
 MaxGen = 30;
 global FunctionTables;
 global FunctionStructure;
 global FunctionAccCard;
 global SelectedObjectives;
 \% The circular structure is created
 FunctionStructure =  CreateListFactorsCircularNK(n,4);
 \% Values are read from a file
 FunctionTables = ReadFunctionsFromData('testNK\_fnt\_N50\_k4Inst\_1.txt',
 FunctionStructure,Card);
 \% Auxiliary structure for  evaluation
 FunctionAccCard = FindListCard(FunctionStructure,Card);
 SelectedObjectives = [1:4:48]; \% Only some of the objectives are evaluated
\% General function that uses the global variables.
 F = 'PartialEvaluateGeneralFunction';
 selparams(1:2) =\{0.5,'ParetoRank\_ordering'\};
 edaparams\{1\} = \{'selection\_method','truncation\_selection',selparams\};
 edaparams\{2\} = \{'replacement\_method','best\_elitism',\{'ParetoRank\_ordering'\}\};
 edaparams\{3\} = \{'stop\_cond\_method','max\_gen',\{MaxGen\}\};
 [AllStat,Cache]=RunEDA(PopSize,n,F,Card,cache,edaparams); \end{programf}
\label{ex:MULTOBJT}  \end{example}



 The code shown above illustrates the way in which multi-objective decomposable functions
 can be  generated and used in MATEDA-2.0. In this example, the
 multi-objective function is based on the NK circular fitness landscape \cite{Kauffman:1993}. Each objective
 corresponds to a subfunction where each
variable depends on its $\frac{k}{2}$  previous and $\frac{k}{2}$
subsequent neighbors. The structure is generated using
\texttt{CreateListFactorsCircularNK} method, which together with
other function generator methods is explained in the next section.
Only a subset of $12$ objectives is selected for optimization. The
default probabilistic model (Tree-EDA) is used with
\texttt{ParetoRank\_ordering} selection.



\subsubsection{Function generator methods}

 Function generator methods can be defined by the user to generate
 multiple instances of a given decomposable function. MATEDA-2.0
 includes a number of auxiliary methods with this purpose. These
 methods can also be used to save and read the structure and values
 of the function from files. A multi-objective implementation of the NK landscape model is
 included as an example of how to use the function generator
 methods.


 \begin{itemize}
 \item \texttt{CreateNKFunctions}:   Generates random values for the table entries of the NK
 landscape.
 \item \texttt{CreateGaussianValuesForFactors}: Generate function values sampling from a Normal
 distribution.
 \item \texttt{CreateListFactorsNK}:   Creates the structure of an instance of the NK random
 landscape, i.e. the $k$ neighbors for each of the variables are  randomly
 selected.
 \item \texttt{CreateListFactorsCircularNK}:   Creates the structure of an instance of the NK landscape
                                where each variable depends on its
                                $\frac{k}{2}$  previous and $\frac{k}{2}$ subsequent neighbors in a
                                circular way.
 \item \texttt{CreateRandomFunctions}: Generates uniform random values for the table entries of a given
 structure.
 \item \texttt{SaveFunctionStructure}:   Saves the structure of a function in the form of a factor
 graph \cite{Kschischang_and_Frey:1998}.
 \item   \texttt{AllNKInstances}:   Generates $m$ instances of a random NK model.  The structure and function of all instances are saved  in
   files.
  \item \texttt{ReadFunctionsFromData}:    Reads the  values  of a  function
  from a file.
\item \texttt{ReadFactorGraphFromData}:    Reads the  structure  of a
given function from a file.
\end{itemize}


\section{How to use MATEDA-2.0 for a given problem}


\subsection{Steps to run an EDA in MATEDA-2.0}

In order to solve a problem using MATEDA-2.0, the following steps
should be followed:

\begin{enumerate}
  \item Create or define the file of the function to be optimized.
  \item Define the type of representation to be used.
  \item Create the vector with the range of values (for the continuous case)
  or the cardinality of the variables (for the discrete case).
  \item Choose each EDA component and determine the
  parameters to be passed to the algorithm.
  \item Execute RunEDA.m.
\end{enumerate}


 MATEDA-2.0 does function maximization. For minimization problems, the fitness function $\hat{f}({\bf{x}})$ has to be modified (e.g. $f({\bf{x}}) = -\hat{f}({\bf{x}})$).
The choice of the EDA components depends on several factors. Among them:
\begin{itemize}
 \item Type of variable representation (discrete or continuous).
 \item Domain
of definition for each variable (discrete problems of high
cardinality may not be treated using complex models).
 \item Existence
of prior information about the problem (e.g. when a feasible
factorization is known it can be employed).
 \item Computational cost of the fitness function.
 \end{itemize}


\subsection{Examples of the optimization problems implemented}


This section presents some examples of the application of MATEDA-2.0
to a number of optimization problems. Some of the problems
implemented and available with the code are the following:

\begin{itemize}
  \item \texttt{EvalIsing}: Ising model \cite{Ising:1925}.
  \item \texttt{EvaluateEnergy}: HP simplified protein model \cite{Dill:1985}.
  \item \texttt{EvaluateEnergyFunctional}: Functional protein model \cite{Hirst:1999}.
  \item \texttt{GaussProtein}: Off-lattice HP protein model \cite{Stillinger_et_al:1993}.
  \item \texttt{EvaluateSAT}: Multi-objective max 3-satisfiability problem \cite{Selman_et_al:1992}.
 \end{itemize}


In the next sections we show different EDA approaches to some of
these problems.


\subsubsection{Tree-EDA for the Ising model}


 The generalized Ising model \cite{Ising:1925} is described by the energy functional
(Hamiltonian):

 \begin{equation}
 E = - \sum_{i<j \in L} J_{ij} \sigma_i \sigma_j - \sum_{i \in L} h_{i}
 \sigma_i  \label{ISINGM}
 \end{equation}
where $L$ is the set of sites called a lattice. Each spin variable
$\sigma_i$ at site $i \in L$ either takes the value $1$ or $-1$. One
specific choice of values for the spin variables is called a
configuration. The constants $J_{ij}$ are the interaction
coefficients.  The ground state is the configuration with minimum
energy. We address the problem of maximizing $f(x) = -E(x)$. The
Ising model has been extensively employed as a testbed of EDAs
\cite{Brownlee_et_al:2008,Hoens_et_al:2007,Mendiburu_et_al:2007a,Muehlenbein_and_Hoens:2005,Pelikan_and_Hartmann:2006,Shakya_and_McCall:2007}.

Example~\ref{ex:Ising} describes the application of the Tree-EDA to
the optimization of the Ising model. The algorithm uses the most
probable configuration sampling \cite{Mendiburu_et_al:2007a}  (i.e.
  the most probable individual from the net is added to the population
  during sampling). The algorithm stops when the optimum is found or the
 maximum number of generations is reached.


\begin{example}[Tree-EDA for the Ising problem]

 \begin{programf}
 global lattice; \% The lattice and the interactions
 global inter;   \% are defined as global variables
 Pop-Size = 500;
 n = 64;
 cache  = [0,0,0,0,0];
 Card = 2*ones(1,n);
 F = 'EvalIsing';  \% Ising function
 MaxGen = 150;
 MaxVal = 86;
 [lattice, inter] = LoadIsing(n, 1);     \% The Ising instances is read
 stop\_cond\_params = \{MaxGen,MaxVal\}; \% The algorithm stops when
                                 \% optimum is found or maxgen is reached \UnnumLine  
 edaparams\{1\} = \{'stop\_cond\_method','maxgen\_maxval',stop\_cond\_params\};
 edaparams\{2\} = \{'sampling\_method',' SampleMPE\_BN',\{PopSize\}\};
 [AllStat,Cache]=RunEDA(PopSize,n,F,Card,cache,edaparams); \end{programf}

\label{ex:Ising} \end{example}

The code above implements the application of the Tree-EDA to the
optimization of the Ising model. A characteristic of this
application, and a suggested solution for problems where some extra
information is needed during the function evaluation step, is the
use of global variables. For the Ising problem implementation,
variables \emph{$lattice$} and \emph{$inter$} are accessed while
reading the Ising instance from a file and during the evaluation
step.


\subsubsection{FDA optimization of the Hydrophobic-Polar (HP) protein model}


We present in this section an FDA application to the HP protein
folding problem \cite{Dill:1985}. It can be approached as
 a discrete problem with constraints and we use a Markov-chain-like probabilistic model
as the structure of the FDA
\cite{Santana_et_al:2004,Santana_et_al:2008a}.


 The HP protein model \cite{Dill:1985} considers two types of
residues: hydrophobic (H)  residues  and   hydrophilic or polar (P)
residues. In this model, a protein is considered a sequence of   these two types of
residues, which are located in regular lattice models forming
self-avoided paths. Given a pair of residues, they are considered
neighbors if they are adjacent  either in the chain  (connected
neighbors) or  in the lattice but  not connected in the chain
(topological neighbors).

In the optimization approach, the search for the protein structure
is transformed into the search for the optimal configuration given
an energy function.  For the HP model, an energy function that
measures the interaction
  between topological  neighbor residues is defined  as  $\epsilon_{HH}=-1$ and
  $\epsilon_{HP}=\epsilon_{PP}=0$.

  The HP problem consists of
  finding the solution that minimizes the total energy. In the linear representation of the sequence,
  hydrophobic residues are represented with the letter H and polar ones, with P. In the graphical representation, hydrophobic proteins are
  represented  by black beads and polar proteins, by white
  beads.

In our representation, for a
 given sequence and lattice,  $X_i$ will represent
 the relative move of residue $i$ in relation to the previous two
 residues.  Taking as a reference the location of the previous two
 residues in the lattice, $X_i$ takes values in $\{0,1, \ldots,
 z-2\}$, where $z-1$ is the number of movements allowed in  the given
 lattice. These
 values respectively mean that the new residue will be located
 in one of the $z-1$ numbers of possible directions with respect to the
 previous two locations.  A  solution  ${\bf{x}}$ can be seen as  a walk in the lattice,
 representing one possible folding of the protein. The codification
 used is called relative encoding, and has been experimentally compared to absolute
 encoding in \cite{Krasnogor_et_al:1999}, showing better results.

Folder $../MATEDA/functions/protein$ contains an implementation of
HP protein model that can be used with different EDA
implementations.


\begin{figure*}
  \begin{center}
   {{\includegraphics[width=8.0cm]{HPFigEx.eps}}}
      \caption{Best solution found by the Markov chain FDA after $100$ generations.}
     \label{fig:HPEx}
    \end{center}
\end{figure*}

Example~\ref{ex:HPMKEDA} shows the code that implements the Markov
chain FDA for an HP sequence of $64$ residues. Other sequence
configurations can be tried by modifying $InitConf$ variable. The
cardinality of the variables is $3$ and the Markov model considers
dependence of each variable to the previous one. The fixed structure
of the model ($Cliques$) is constructed using  the
\texttt{CreateMarkovModel} method.

 The EDA uses a repairing method based on backtracking \cite{Cotta:2003} which tries to enforce
 that each sequence  is folded forming a self-avoided path in the
 lattice. At the end of the run, the solution found by the EDA is
 visualized using the method \texttt{PrintProtein}. Figure~\ref{fig:HPEx} shows the
 best solution found by one run of this algorithm after $100$
 generations.

\begin{example}[Markov Chain FDA for the HP protein model]
\begin{programf}
 global InitConf;   \% This is the HP protein instance
                    \% defined as a sequence of zeros and ones
 InitConf =   [zeros(1,12),1,0,1,0,1,1,0,0,1,1,0,0,1,1,0,1,1,0,0,1,1
 \UnnumLine   ,0,0,1,1,0,1,1,0,0,1,1,0,0,1,1,0,1,0,1,zeros(1,12)];
 PopSize = 500;
 NumbVar = size(InitConf,2);
 cache  = [0,0,0,0,0];
 Card = 3*ones(1,NumbVar);
 maxgen = 100;
 \% The Markov chain model(Cliques) is constructed specifying the number of
 \% conditioned (previous) variables. In the example below this number is
 \% 1., i.e. $p(x) = p(x0)p(x1|x0) ... p(xn|xn-1)$
 Cliques = CreateMarkovModel(NumbVar,1);
 F = 'EvaluateEnergy'; \% HP protein evaluation function

 edaparams\{1\} = \{'learning\_method','LearnFDA',{Cliques}\};
 edaparams\{2\} = \{'sampling\_method','SampleFDA',{PopSize}\};
 \% Repairing method used to guarantee that solutions do not self-intersect
 edaparams\{3\} = \{'repairing\_method','HP\_repairing',\{\}\};
 edaparams\{4\} = \{'stop\_cond\_method','max\_gen',\{maxgen\}\};
[AllStat,Cache]=RunEDA(PopSize,NumbVar,F,Card,cache,edaparams)

 \% To draw the resulting solution use function PrintProtein(vector),
 \% where vector is the best solution found.
   vector = AllStat\{maxgen,2\};
   PrintProtein(vector); \end{programf}

 \label{ex:HPMKEDA} \end{example}

\subsubsection{Bayesian network based EDA for a multi-objective 3-satisfiability problem}

Let ${\bf{U}} = \{U_1,U_2 \cdots U_n\}$ be a set of $n$ Boolean
variables. A (partial) truth assignment for ${\bf{U}}$ is a
(partial) function $T: {\bf{U}} \rightarrow \{true,false\}$.
Corresponding to each variable $U_i$ are two literals, $u_i$ and
$\lnot u_i$. A literal $u_i$ (resp. $\lnot u_i$) is $true$ under $T$
iff $T(u_i)=true$ (resp. $T(u_i)=false$). A set of literals is
called a clause, and a set or sequence (tuple) of clauses a formula.

Let $\phi $ be a formula, ${\bf{U}} = vars(\phi )$, and $C$ a clause
in $\phi $. We interpret $\phi $ as a formula of the propositional
calculus in conjunctive normal form, so that a truth assignment $T$
for ${\bf{u}}$ satisfies $C$ iff at least one literal $u_i \in C$ is
true under $T$, and $T$ satisfies $\phi $ iff it satisfies every
clause in $\phi $. The satisfiability problem (SAT)
\cite{Selman_et_al:1992}  is the problem of finding a satisfying
assignment for a formula.

In the multi-objective version of the SAT problem, there are $m$
formulas $(\phi_1, \dots, \phi_m)$, each representing a different objective. This type of
problems have been employed to test EDAs for multi-objective
problems \cite{Bosman:2003}. MATEDA-2.0 includes the method
\texttt{EvaluateSAT} which evaluates a solution on a set of 3-SAT
formulas contained in the global variable $Formulas$. The output is
a multi-objective solution, one component corresponds to the
evaluation of one formula. The method \texttt{MakeRandomFormulas}
can be used to generate random instances of the multi-objective SAT
problem and a number of these instances can be found in the folder
$../MATEDA/functions/SAT/instances$.

Example~\ref{ex:SATEDA} shows the code that implements a Bayesian
network based EDA  for multi-objective 3-SAT. The global variable
$Formulas$ stores the formulas that are read from the file. In this
example, there are $20$ binary variables, $10$ formulas and $20$
clauses in each formula. ParetoRank\_ordering is used as the
ordering criterion and truncation selection is applied.

\begin{example}[Bayesian network based EDA  for multi-objective 3-SAT]
\begin{programf}
 global Formulas;   \% Global variable for SAT function
 load TypeForm\_1\_Form\_1.mat; \% Multimodal SAT instance
 n = 20;  \% Number of variables
 m = 10;  \% Number of objectives
 c = 20;  \% Number of clauses
 PopSize = 500;
 NumbVar = n;
 cache  = [1,1,1,1,1];
 Card = 2*ones(1,NumbVar);
 maxgen = 50;
 F = 'EvaluateSAT'; \% 3-SAT function
 selparams(1:2) = \{0.5,'ParetoRank\_ordering'\};
 sampling\_params(1:3) = \{PopSize,m,Obj\_Card\};
 BN\_params(1:7) = \{'k2',5,0.05,'pearson','bayesian','no'\};
 edaparams\{1\} = \{'learning\_method','LearnBN',BN\_params\};
 edaparams\{2\} = \{'sampling\_method','SampleBN',sampling\_params\};
 edaparams\{3\} = \{'selection\_method','truncation\_selection',selparams\};
 edaparams\{4\}=\{'replacement\_method','best\_elitism',\{'ParetoRank\_ordering'\}\};
 edaparams\{5\} = \{'stop\_cond\_method','max\_gen',\{maxgen\}\};
 [AllStat,Cache]=RunEDA(PopSize,NumbVar,F,Card,cache,edaparams); \end{programf}
\label{ex:SATEDA} \end{example}



\subsubsection{Mixture of multivariate Gaussian distributions for a spacecraft trajectory optimization problem}

Spacecraft trajectory problems can be posed as global optimization problems with constraints.
One class of these problems is the multiple gravity assist missions with the possibility of using
deep space manoeuvres (MSGDSM) \cite{Vinko_et_al:2007}. It consists of finding an interplanetary
trajectory of a spacecraft equipped with chemical
propulsion, able to thrust its engine once at any time between each
trajectory leg.

The generic form of the MGADSM
problem can be written as \cite{Vinko_et_al:2007}:

\begin{equation}
\begin{tabular}{cc}
find & $x \in \Re^n$  \\
minimize & $J(x)$ \\
subject to & $g(x)$ \\
\end{tabular}  \nonumber
\end{equation}
where $x$ is the decision vector, $J$ is the objective function and $g$ are non linear
constraints that may come from operational considerations or from the spacecraft system
design. These problems have many local optima and
 constitute a formidable benchmark for the evaluation of
continuous EDAs and other evolutionary algorithms.

We choose an instance of this problem corresponding to the design of a deltaV-EGA manoeuvre
required to reach Jupiter using an Earth Earth Jupiter fly-by sequence with deep space
manoeuvres. An EDA based on the use of a mixture of multivariate Gaussian distributions is selected
to address the problem.

Example~\ref{ex:TRAJEDA} shows the code of the EDA. First, a global variable that will store
the information about the problem instance is defined. Then, a file containing the problem
 instance is loaded\footnote{This instance and the  Matlab method that
 allows the evaluation of the function can be downloaded from
 www.esa.int/gsp/ACT/inf/op/globopt.htm}. The problem has $12$ continuous variables
 with ranges defined in $Card$. The learning parameters of the EDA model specify that:
 Points are clustered according to the decision variables values,
 the clustering method used to learn the mixture is
 k-means,  the number of mixture components is $10$ and the measure used to cluster the points
 is the $sqEuclidean$ metric as defined in Matlab.
 Once the mixture of multivariate Gaussian distributions have been learned and sampled, solutions
 are repaired to guarantee the variables values are within the prescribed bounds.



\begin{example}[Mixture of Gaussian distributions for a trajectory problem]
\begin{programf}
 global MGADSMproblem \% Global variable for the space trajectory problem
 load EdEdJ;    \% The instance is read
 NumbVar = 12;
 PopSize = 5000;
 F = 'EvalSaga';
 Card(1,:) = [7000,0,0,0,50,300,0.01,0.01,1.05,8,-1*pi,-1*pi];
 Card(2,:) = [9100,7,1,1,2000,2000,0.90,0.90,7.00,500,pi,pi];
 cache  = [0,0,1,0,1];
 learning\_params(1:5) = \{'vars','ClusterPointsKmeans',10,'sqEuclidean',1\};
 edaparams\{1\} = \{'learning\_method','LearnMixtureofFullGaussianModels',learning\_params\};
 edaparams\{2\} = \{'sampling\_method','SampleMixtureofFullGaussianModels',\{PopSize,3\}\};
 edaparams\{3\} = \{'replacement\_method','best\_elitism',\{'fitness\_ordering'\}\};
 selparams(1:2) = \{0.1,'fitness\_ordering'\};
 edaparams\{4\} = \{'selection\_method','truncation\_selection',selparams\};
 edaparams\{5\} = \{'repairing\_method','SetWithinBounds\_repairing',\{\}\};
 edaparams\{6\} = \{'stop\_cond\_method','max\_gen',\{5000\}\};
 [AllStat,Cache]=RunEDA(PopSize,NumbVar,F,Card,cache,edaparams)
\end{programf}
 \label{ex:TRAJEDA} \end{example}



%\section{Analysis and visualization of the data generated in EDAs}



\section{Analysis of data related to the points generated and the evaluation
of the (possible multiple) objective functions}


 In this section, we present a number of methods used to extract, process, and visualize the information gathered during the EDA evolution. This information can reveal  new knowledge about the problem structure and may serve to evaluate the efficacy of the EDA approach and the role played by its different components.

 Visualization techniques are an important tool for interpretation
 and understanding of data. Classical approaches dealing with
 visualization of small, isolated problems have recently shifted to
 the visualization of massive scale, dynamic data, comprised of
 elements of varying levels of certainty and abstraction
 \cite{Buerger_and_Hauser:2007,Thomas_and_Cook:2006}.

 It is a complex task to find a clear separation between the different
classes  of data generated during the EDA search. Nevertheless, we
use the following classification to distinguish the way the
different classes of data are produced:

\begin{itemize}
  \item Data related to the points generated and the evaluation of the (possible multiple)
  objective functions (e.g. number and distribution of the points generated, shape of the
  Pareto front approximations, etc.).
 % \item Statistics describing the components of the algorithm  (e.g. number of new solutions added to the selected set,time spent by each method, etc.).
  \item Probabilistic models (e.g. structure and parameters of the probabilistic models).
\end{itemize}

 From the previous classification, it is possible to implement
different methods to extract useful information from the EDA search
and  visualize this information. Methods related to the first
class are analyzed in this section, methods to analyze and visualize
the structures and parameters of the probabilistic models are
presented in the next section.


The main classes of methods for the analysis of  data related to
the points generated and the evaluation of the objective functions are the following:

   \begin{itemize}
     \item Computation and visualization of fitness related measures.
     %\item Computation and visualization of measures of convergence.
     \item Analysis and visualization of dependencies between variables and correlations between the
     objectives.
   \end{itemize}


\subsection{Fitness related measures}


Fitness related measures are obtained from the fitness values of the
solutions visited by the algorithm at each generation.

The  average fitness ($\bar{f}$) of the population at each
generation can be used as a source of information about the behavior
of the EDA. If we are looking for a maximum, an increase in the average fitness means that the
algorithm is able to generate better solutions. However, an increase
of the average fitness can hide a loss of diversity in the
population. In this case, the fitness variance ($\sigma(f)$) can
support additional information about whether the fitness values of
the population are really diverse. Similarly, the distribution of
the fitness function represented using histograms gives a clearer
perspective of the diversity in the population.

The response to selection  ($R(t)$) is a general measure of the
improvements obtained in the average fitness of the population by
the application of evolutionary algorithm operators. The amount of selection allows
to measure the effect of the selection operator in terms of the
fitness ($S(t)$). The realized heritability ($b(t)$) can serve to
evaluate the effect of the sampling and replacement methods. The
mathematical framework that involves the use of $R(t)$, $S(t)$ and
$b(t)$ was originally proposed in population genetics and has been
applied  to the analysis of the breeder genetic algorithm
\cite{Muhlenbein_and_Voosen:1993a} and EDAs
\cite{Muhlenbein:1997,Santana:1999,Santana_et_al:2008c}.



The following methods that compute fitness related measures are
included in MATEDA-2.0:

   \begin{itemize}
   \item \texttt{Mean\_Var\_Objectives}: Computes the average fitness and variance of the objectives
    ($\bar{f}$,$\sigma(f)$).
    \item \texttt{ObjectiveDistribution}: Computes and visualizes the fitness distribution
    for a given range of generations and a subset of objectives using histograms.
    \item \texttt{Response\_to\_selection}: Computes the response to selection for every objective, $R(t) = \bar{f}(t+1)
    -\bar{f}(t)$.
    \item \texttt{Amount\_of\_selection}: Computes the amount of selection for every objective, $S(t) = \bar{f}_{s}(t)
    -\bar{f}(t+1)$.
    \item \texttt{Mean\_Var\_Objectives}: Computes the realized heritability for every objective, $b(t) =
    \frac{R(t)}{S(t)}$.
   \end{itemize}



Example~\ref{ex:COMPFIT}  shows the MATEDA-2.0 code for the
computation of fitness related measures from the data generated by a
Tree-EDA (the EDA that MATEDA-2.0 uses when no learning and sampling methods are specified).

\begin{example}[Computation of fitness related measures]
\begin{programf}
  PopSize = 300;
  n = 45;
  cache  = [0,0,0,1,1];
  Card = 2*ones(1,n);
  F = 'sum'; \% Onemax function;
  ngen = 10;
  edaparams\{1\} = \{'stop\_cond\_method','max\_gen',\{ngen\}\};
  [AllStat,Cache]=RunEDA(PopSize,n,F,Card,cache,edaparams)
   for i=1:ngen,
     auxr\{1,i\} = Cache\{4,i\};
     auxs\{1,i\} = Cache\{5,i\};
   end,
  [mean\_obj,var\_obj] = Mean\_Var\_Objectives(auxr);
  [RS] = Response\_to\_selection(auxr);
  [S] = Amount\_of\_selection(auxr,auxs);
  [b] = Realized\_heritability(auxr,auxs);
  ObjectiveDistribution(auxr,1,[1:ngen]);
\end{programf}
\label{ex:COMPFIT} \end{example}

The code presented in Example~\ref{ex:COMPFIT} starts by executing a UMDA with a maximum
of $10$ generations for the $Onemax$ function. Since $cache  =
[0,0,0,1,1]$, the evaluation of the entire and selected populations
for all generations are stored. From this information, auxiliary
variables $auxr$ and $auxs$ are computed and from them the fitness
related measures have been calculated.


\subsection{\emph{A posteriori} analysis of the dependencies}

 \emph{A posteriori} modeling of the dependencies that arise in the
 selected sets (or the final Pareto sets approximations found by EDAs), can be useful
 to study the efficiency of EDAs that use different classes of models, to
 construct concise descriptions of the set of optimal solutions (e.g. graphical
 models of the obtained Pareto set approximations), to extract and visualize
characteristic features of the selected sets at each
generation, etc.  We illustrate
 this claim in this section and explain the way in which MATEDA-2.0
 allows the user to do \emph{a posteriori} modeling of the dependencies

The most straightforward approach to modeling the dependencies
between the variables is to construct a  probabilistic graphical
model, as usually done in \emph{online} modeling employed by EDAs.
However, a probabilistic graphical model is not the only choice
available to model the dependencies in the data. For instance, a
matrix of correlations between the objectives has been employed
\cite{Deb_and_Saxena:2005,Lopez_et_al:2008} in multi-objective
optimization to detect conflicting and redundant objectives.

Notice also, that in contrast to the \emph{online} modeling case, in
\emph{a posteriori} modeling we do not require to model the selected
sets learned at all the generations, and in general the main
objective of modeling is not to sample new solutions.

 The following methods have been implemented in MATEDA-2.0 to model (\emph{a
 posteriori}) dependencies.

\begin{itemize}
\item \texttt{Corr}: This Matlab function computes the correlation matrix between a
set of variables given a set of observations:
\item Methods for learning probabilistic models: Currently, MATEDA-2.0 uses the capabilities of the Bayesian
 networks toolbox for representing dependencies for continuous and mixed
 variables.  The following MATEDA-2.0 functions can be used to learn the models:
\begin{itemize}
  \item \texttt{LearnBN}.
  \item \texttt{LearnGaussianNetwork}.
  \item \texttt{LearnMixedModel}.
\end{itemize}
\item \texttt{ShowParallelCoord}: Parallel coordinate visualization of the
objectives.
\end{itemize}

The correlation matrix of the Pareto set approximation obtained with
function \texttt{Corr} can provide an initial idea of which
objectives are contradicting and redundant. However, extracting a
set of non-redundant objectives from this matrix is not
straightforward. Different techniques have been proposed in
multi-objective optimization with this purpose  \cite{Brockhoff_and_Zitzler:2006,Deb_and_Saxena:2005,Lopez_et_al:2008,Purshouse_and_Fleming:2003}.

The methods used for learning the models have been analyzed in
Section~\ref{sec:LEARNMETHODS}. Its application to \emph{a
posteriori} modeling is very similar but the user should take into
account the difference between the problem objectives and the
problem variables in the analysis of the model. Also, the type of
representation used by objectives and variables have to be
considered in the selection of the probabilistic model.

\subsubsection{Parallel coordinate visualization of objectives}
\label{sec:ParCoord}

In parallel coordinates \cite{Inselberg:2009}, every observation is
plotted for each axis/variable,
  and a connecting line is drawn for each observation between all the
  axes. Parallel coordinates can be very effective to detect
  redundant and conflicting objectives. It can also be used to identify
  outlier instances. For example, those sets of objective values
  that do not follow the same trend than the rest.

    An important issue in the effectiveness of the parallel coordinates approach to the visualization
  of multivariate data is the order and arrangement of dimensions (which in our case corresponds to the objectives). Ordering may help to reduce cluttering,
improving  visualization.   One way to deal with this problem is to
rearrange the data dimensions such that dimensions showing a similar
   behavior are positioned next to each other. However, the problem of finding an optimal one- or two-dimensional
     arrangement of the dimensions based on their similarity is NP-complete \cite{Ankerst_et_al:1998}.
Heuristic algorithms have been proposed to find satisfying solutions
\cite{Ankerst_et_al:1998}.

In MATEDA-2.0, the method for displaying the parallel coordinates of
the objectives is invoked as:

 \texttt{ShowParallelCoord(fs,AllObjectives(:,ordering));}
where $fs$ is the font size used in the parallel coordinates
representation and $ordering$ is an order of the objectives.

MATEDA-2.0 includes a function that allows the user to define a
method to find an  order of the objectives. The
\texttt{ClusterUsingDist} method clusters together objectives with
strong similarity in the Pareto set approximation (or selected
population). Clustering is done according to a parameter $distance$
which can take as values  any of the distances used by Matlab
command \texttt{pdist} (ex. 'correlation', 'euclidean', etc.).
\texttt{ClusterUsingDist} uses affinity propagation
\cite{Frey_and_Dueck:2007}  to cluster the objectives.

Example~\ref{ex:MULTPARETO} shows how to find a Pareto set
approximation using the best solutions found by an EDA. The code
shown is the continuation of the code included  in
example~\ref{ex:MULTOBJT} where the generation and optimization of a
multi-objective decomposable function is presented.

\begin{example}[Construction of the Pareto set approximation]
 \begin{programf}
  AllSols = [];
  AllVals = [];
 for i=1:MaxGen,
  AllSols = [AllSols;Cache\{1,i\}];   \% All the populations
  AllVals = [AllVals;Cache\{4,i\}];   \% All the evaluations
 end

 \% The set of Pareto solutions found by the EDA is extracted
 Index = FindParetoSet(AllSols,AllVals);
 ParetoPop = AllSols(Index,:);
 ParetoVals = AllVals(Index,:);
 \% The Pareto set approximation is shown using parallel coordinates
 parallelcoords(ParetoVals);
  \% The correlations between the objectives are computed.
 ObjectivesCorr = corr(ParetoVals);
  \% Affinity propagation is used to reduce the number of objectives
  \% by selecting exemplars of the correlated variables
 [idx,netsim,dpsim,expref]=apcluster(ObjectivesCorr,mean(ObjectivesCorr)); \end{programf}
 \label{ex:MULTPARETO} \end{example}


 The code shown above illustrates the way in which  a Pareto set approximation
 can be obtained from the selected solutions found in every generation of an EDA.
 In this example, the
 multi-objective function is based on the NK circular fitness landscape. Each objective
 corresponds to a subfunction where each
variable depends on its $\frac{k}{2}$  previous and $\frac{k}{2}$
subsequent neighbors.

 After the Pareto set approximation has been found, it is visualized using  parallel coordinates
 and the correlation between its objectives are found. Finally, affinity propagation is applied
 to reduce 'redundant' objectives.

%\section{Statistics describing the components of the algorithm}

\section{Analysis of the probabilistic models}

Work on the analysis  of the probabilistic models learned by EDAs
and its relationship with the problem structure for single objective
functions, have been accomplished for directed
\cite{Bengoetxea:2003,Echegoyen_et_al:2007,Echegoyen_et_al:2008,Hauschild_and_Pelikan:2008,Hauschild_et_al:2008,Lima_et_al:2007,Muehlenbein_and_Hoens:2005},
undirected
\cite{Santana:2005,Santana_et_al:2007b,Santana_et_al:2008} and
Gaussian \cite{Bengoetxea:2003} models. This work has been mainly
focused on the number and frequency of the dependencies represented
in the graphical structure of the model. Recent work explores other descriptors of the learned
models such as the correlation between the model prediction and the fitness \cite{Brownlee_et_al:2008}, the likehood of the model \cite{Lima_et_al:2008},
and the probabilities given by the model to some relevant points of the search \cite{Echegoyen_et_al:2009}.  MATEDA-2.0 allows the user to conduct the type of
analysis achieved in previous work but also permits the detection of
frequent substructures in the models learned, the visualization of
interactions between edges in the models and other additional
features.


\subsection{Data structures to represent the models}

  The basic  information about the interactions between the variables of a problem in a given model
  is represented in MATEDA-2.0 using an edge matrix $E$, where  $E(i,j)=1$ and $E(j,i)=1$
  if and only if  variables $X_i$ and $X_j$ satisfy some previously defined  structural relationship.
  Otherwise,  $E(i,j)=0$ and $E(j,i)=0$.

  The structural relationship depends on the type of probabilistic
  model. The following are some examples of relationships:
  \begin{itemize}
    \item Factorized model: Two variables are related if they appear
    together in any of the factors.
    \item Markov network model: Two variables are related if one of
    them belongs to the neighborhood of the other variable.
    \item Bayesian networks: Two variables are related is there is
    an arc between them.
  \end{itemize}

 Notice that, for Bayesian and Gaussian networks,  we disregard the direction of the
 dependencies. This is a common approach in the analysis of this
 type of directed graphical models in EDAs \cite{Echegoyen_et_al:2007,Echegoyen_et_al:2008,Hauschild_and_Pelikan:2008,Hauschild_et_al:2008,Lima_et_al:2007}.

Let us suppose, we have all the models learned in $nruns$ executions of
an EDA with at most $maxgen$ iterations (when using MATEDA-2.0 with
the option $cache(3)=1$, all the models of a single execution are
saved in  the cell array $Cache\{3,:\}$). From this data, the data
structure $run\_structures$ can be  generated.

\begin{itemize}
 \item $run\_structures\{1\} = indexmatrix(n,n)$: Associates an index to each possible edge in the network.
 e.g. $indexmatrix(1,2) = 1$, number of edges $m = \frac{n(n+1)}{2}$.
 \item  $run\_structures\{2\} = AllBigMatrices\{nruns\}(m,maxgen)$:
 For each run contains whether the edge $e$ appeared in generation $j$.
 \item $run\_structures\{3\} = AllSumMatrices(m,maxgen)$, i.e. the number of runs that each edge $e$ appeared in generation
 $j$.
 \item $run\_structures\{4\} = AllContactMatrix\{maxgen\}(n,n)$: The number of runs in which edge
 $i,j$  was present in generation $k$.
 \item $run\_structures\{5\} =  SumAllContactMatrix(n,n)$, i.e. the
 total number of times edge $(i,j)$ was present in all the structures learned in all generations of all runs.
\end{itemize}

The following programs allow to extract from the models learned by
different type of EDAs the information which is stored in the
$run\_structures$.

\begin{itemize}
  \item \texttt{ReadMNStructures}:  Extracts the information from  the Markov networks or factorization models learned
  during the execution of MATEDA-2.0 (i.e. Cliques saved in $Cache\{3,:\}$).
  \item \texttt{ReadBNStructures}:  Extracts the information from  the Bayesian or Gaussian networks learned
  during the execution of MATEDA-2.0 (i.e. the models saved in $Cache\{3,:\}$).
  \item \texttt{ReadStructures}:  Extracts the information from  a file with the edges corresponding
  to the structures learned by an arbitrary EDA.
\end{itemize}

The program \texttt{ReadStructures} is particularly useful when an
implementation of EDAs not contained in MATEDA-2.0  has been applied
and we want to use MATEDA-2.0 to analyze the structures.  In this
case, the structure must be provided in a file that contains two columns of numerical
values, where a row $(i,j)$  means that there is an edge between node $i+1$ and $j+1$ (note that
vertices in the file are numbered from $0$ to $NumberVar-1$). A negative row  $(-g, -r)$
indicates that the previous (positive) rows were learned in generation $g$ of run $r$. Example~\ref{ex:STRUCTFILE} shows a file
containing different structures.

\begin{example}[File of some model structures learned by EDAs]
 \begin{programf}*
             2  0
             8  0
             14 3
             -1 -1
             13 8
             0 10
             3 5
             9 10
            -1 -2 \end{programf}
File that contains the structures learned by a
 Bayesian network at different generations and runs.
 \label{ex:STRUCTFILE} \end{example}

 The structure learned at run $1$, generation $1$ contains edges
          $(3,1)$, $(9,1)$ and $(15,4)$.  The structure learned at run $2$,
          generation 1 contains edges  $(14,9)$, $(1,11)$, $(4,6)$ and $(10,11)$.

\subsection{Methods implemented for the analysis and visualization
of the structures}

The following methods have been implemented:

\begin{itemize}
 \item \texttt{ViewSummStruct}: Shows one image where each edge has a color proportional
  to the times it has been present in the structures learned in all
  generations.
 \item \texttt{ViewInGenStruct}: Shows images where each edge has a color proportional
 (relative to $nruns$)
  to the times it has been present in the structures
  learned in those generations specified by the user. There is one figure for each generation.
 \item \texttt{ViewEdgDepStruct}: Searches for a given substructure in the set of all the
  structures learned,
 and shows the adjacency matrices corresponding to the structures that contain it.
 The substructure is specified by giving  the edges that are present/absent in them.
 Different visualization options are implemented.
 \item \texttt{ViewPCStruct}: Searches for substructures in the set of all the structures learned
 and shows the parallel coordinates of the edges and the generations at which they are learned.
 The minimal number of times that an edge has to appear in (all) the
structures learned and the method used to order the variables before
displaying them using  parallel coordinates are parameters of the
algorithm.
\item \texttt{ViewDenDroStruct}: Shows the dendrograms of the edges according to
                       their co-occurrence in the structures learned by
                       the EDAs. Allows to detect complex hierarchical
                       relationships between the variables of the problem.

\item \texttt{ViewGlyphStruct}: Shows the glyph representation of a subset of edges
learned at a given set of runs and generations.
\end{itemize}

\subsubsection{\texttt{ViewSummStruct} method}

 The type of structure shown by the \texttt{ViewSummStruct} method is usually called frequency
  matrix \cite{Echegoyen_et_al:2007,Echegoyen_et_al:2008,Santana_et_al:2007b} or probabilistic coincidence matrix \cite{Hauschild_et_al:2007}.
 It serves to detect frequent patterns  of interactions in the models learned during
 the search. Example~\ref{ex:FreqMat} describes a case where the
 frequency matrix has been computed from the Bayesian network
 structures learned in several executions of the EBNA-Exact
 algorithm  \cite{Echegoyen_et_al:2007,Echegoyen_et_al:2008}.


\begin{figure*}
  \begin{center}
   {{\includegraphics[width=8.0cm]{FreqMatrixExample.eps}}}
      \caption{Frequency matrix computed from the structures
 learned by EBNA-Exact.}
     \label{fig:FreqMat}
    \end{center}
\end{figure*}

\begin{example}[Application of \texttt{ViewSummStruct} method]
 \begin{programf}
  NumberVar = 20;
  [run\_structures,maxgen,nruns] = ReadStructures('ProteinStructsExR.txt',NumberVar );
  [results] =  ViewStructures(run\_structures,NumberVar ,maxgen,nruns,
  'viewmatrix\_method','ViewSummStruct',{[150,14]}); \UnnumLine \end{programf}
\label{ex:FreqMat} \end{example}

 Figure~\ref{fig:FreqMat} shows the image representation of the frequency matrix computed from a file containing the structures
 learned by EBNA-Exact in all runs, all generations. Bright colors means higher frequencies.


 The frequency matrix representation has two main limitations. First, it is not possible to investigate
 the types of structures learned at each generation. It has been early observed that
 structures of the models learned by EDAs can change along the evolution \cite{Bengoetxea:2003}. The
 other limitation of this type of representation is that
 it is not possible to capture interactions between different
 substructures of the problem. For instance, if the frequency matrix shows that two edges have occurred
 ten times in the structures learned,  we cannot determine how
 often they have appeared together in the same structure.

\subsubsection{\texttt{ViewInGenStruct} method}

 The $\texttt{ViewInGenStruct}$ method addresses the first
 limitation of the frequency matrix representation by  displaying frequency matrices computed from structures learned
 at each generation. From a sequence of images it is then possible
 to distinguish the dynamics of the structures along the evolution.
 Example~\ref{ex:ViewInGens} shows  an example of the application of
 the \texttt{ViewInGenStruct} method.

\begin{figure*}
  \begin{center}
   {{\includegraphics[width=4.0cm]{ViewInGen1.eps}}}
   {{\includegraphics[width=4.0cm]{ViewInGen5.eps}}}
   {{\includegraphics[width=4.0cm]{ViewInGen10.eps}}}
   \caption{Frequency matrices showing the most frequent edges in the
 structures learned by EBNA-Exact at different generations.}
    \label{fig:FreqMatGens}
    \end{center}
\end{figure*}


\begin{example}[Application of \texttt{ViewInGenStruct} method]
 %Frequency matrices showing the most frequent edges in the
 %structures learned by EBNA-Exact at different generations.
 \begin{programf}
   [run\_structures,maxgen,nruns] = ReadStructures('ProteinStructsExR.txt',20);
   [results] = ViewStructures(run\_structures,20,maxgen,nruns, 'viewmatrix\_method',
   'ViewInGenStruct',{[150,14];[1,5,10]});  \UnnumLine \end{programf}
    \label{ex:ViewInGens} \end{example}

 Figure~\ref{fig:FreqMatGens} shows the output of the \texttt{ViewInGenStruct}
 method invoked in Example~\ref{ex:ViewInGens}. The three frequency matrices shown are learned
   in generations $1$, $5$ and $10$ respectively. It can be observed that the number
   of dependencies captured in the models diminish as the
   generations increase.



\subsubsection{\texttt{ViewEdgDepStruct} method}

 The $\texttt{ViewEdgDepStruct}$  serves to detect and display particular substructures learned in
 the models, allowing to identify complex patterns of interactions
 in the problem and addressing the second limitation of the \texttt{ViewSummStruct} method.
 For instance, the user might be interested in testing a particular
 hypothesis related to the mapping \emph{problem structure -
 model structure}. By computing the frequency of appearance of the
  substructures in the models these hypothesis could be
 tested. Example~\ref{ex:ViewEdgDepStruct} shows  the application of
 the \texttt{ViewEdgDepStruct} method.

\begin{figure*}
  \begin{center}
   {{\includegraphics[width=6.0cm]{FreqMatrixSub1Example.eps}}}
   {{\includegraphics[width=6.0cm]{FreqMatrixSub2Example.eps}}}
   {{\includegraphics[width=6.0cm]{FreqMatrixSub3Example.eps}}}
   {{\includegraphics[width=6.0cm]{FreqMatrixSub4Example.eps}}}
   \caption{Edge matrices of those structures learned in several runs
   EBNA-Exact and
 such that edges $(3,4)$ and $(4,5)$ appear together in the structure and edge $(3,5)$ does not
 appear.}
    \label{fig:EdgeMat}
    \end{center}
\end{figure*}


 \begin{example}[Application of \texttt{ViewEdgDepStruct} method]
  \begin{programf}
   viewparams\{1\} = [100,14];
   viewparams\{2\} = [3 4 1; 4 5 1; 3 5 0]; \% The substructure is described
   viewparams\{3\} = [1:nruns]; \% Selected runs (All)
   viewparams\{4\} = [1:maxgen]; \% Selected generations (All)
   viewparams\{5\} = 'all\_graphs'; \% Graphs to be visualized (All)
   [run\_structures,maxgen,nruns] = ReadStructures('ProteinStructsExR.txt',20);
   [results] = ViewStructures(run\_structures,20,maxgen,nruns, 'viewmatrix\_method',
   'ViewEdgDepStruct',viewparams);  \UnnumLine \end{programf}
\label{ex:ViewEdgDepStruct} \end{example}

 Figure~\ref{fig:EdgeMat} shows the four adjacency matrices of those structures
 learned in several runs
  of EBNA-Exact and that satisfy the condition that
 edges $(3,4)$ and $(4,5)$ appear together in the structure and edge $(3,5)$ does not
 appear. The output variable \emph{$results$} describes at
  which run and generation each of the visualized structures have  been found.

\subsubsection{\texttt{ViewPCStruct} method}

The parallel coordinates, explained in Section~\ref{sec:ParCoord},
can also be used to visualize the structure of the models. Method
\texttt{ViewPCStruct} implements a parallel coordinate visualization
in which the   vertical axis represents the generation at which
edges (shown in the horizontal axis) have been learned. A line
between two points means that both edges appear in the same
structure learned at the same generation.

The method allows to select the most relevant subset of edges for
representation. Parameter $const_{edg}$ is the  minimal number of
times that an edge has to appear in (all) the structures learned to
be selected for visualization. Since the  clarity of the parallel
coordinate  visualization depends on the number of variables, this is
an important parameter. Also the minimal number of edges
($min_{edg}>0$) in the substructures selected is a parameter of the
algorithm. Finally, the user can specify the method to order
the variables before displaying them. \texttt{ViewPCStruct} outputs
each of the represented edges together with the generation and run
at which it has  been learned.

Example~\ref{ex:ViewPCStruct} shows an example of the application of
the \texttt{ViewPCStruct} method.

\begin{figure*}
  \begin{center}
   {{\includegraphics[width=8.0cm]{PCedges.eps}}}
     \caption{Parallel coordinate visualization of the most frequent edges in the structures
     learned by EBNA-Exact.}
    \label{fig:PCVis}
    \end{center}
\end{figure*}

\begin{example}[Application of \texttt{ViewPCStruct} method]

 \begin{programf}
 viewparams\{1\} = [14]; \% Font size
 viewparams\{2\} = []; \% The edges will be found by the algorithm
 viewparams\{3\} = 60; \% Those edges that appear at least $60$ times
 viewparams\{4\} = 2;  \% Structures that have at least two edges
 viewparams\{5\} = 'ClusterUsingDist'; \% Variables are clustered using affinity prop.
 viewparams\{6\} = 'correlation'; \% Similarity measure used is the correlation
 [run\_structures,maxgen,nruns] = ReadStructures('ProteinStructsExR.txt',20);
 [results] = ViewStructures(run\_structures,20,maxgen,nruns,'viewmatrix\_method',
 'ViewPCStruct',viewparams);  \UnnumLine \end{programf}
 \label{ex:ViewPCStruct} \end{example}

 Figure~\ref{fig:PCVis} shows the  parallel coordinate visualization of the most frequent edges (they appear
  at least $20$ times) in
   the structures  learned in several runs of EBNA-Exact
     and that satisfy the condition that the structures have at least
     two edges.

\subsubsection{\texttt{ViewDenDroStruct} method}

 Dendrograms are graphs that serves to represent hierarchical trees.
     A dendrogram consists of many U-shaped lines connecting objects in
     the  hierarchical tree  \cite{Johnson:1967}. The height of each U represents the distance between
    the two objects being connected.

Similarly to the \texttt{ViewPCStruct} method,
\texttt{ViewDenDroStruct} allows to select the most relevant subset
of edges for representation. Parameter $const_{edg}$ is the  minimal
number of times that an edge has to appear in (all) the structures
learned to be selected for visualization. Also the minimal number of
edges ($min_{edg}>0$) in the substructures selected is a parameter of
the algorithm.

The \texttt{ViewDenDroStruct} method computes first a hierarchical
tree of the selected edges based on a distance defined by the user
(usually the correlation of the edges in the structures learned).
Then, the hierarchical tree is displayed using a dendrogram. This
representation can be very effective to detect hierarchical
relationships between substructures in the models.
Example~\ref{ex:ViewDendroStruct} shows an example of the
application of the \texttt{ViewDendroStruct} method.


\begin{figure*}
  \begin{center}
   {{\includegraphics[width=8.0cm]{DendroViewExp.eps}}}
     \caption{Dendrogram visualization of the most frequent edges in the structures
     learned by EBNA-Exact.}
    \label{fig:DendroVis}
    \end{center}
\end{figure*}


\begin{example}[Application of \texttt{ViewDendroStruct} method]

 \begin{programf}
 viewparams\{1\} = [14]; \% Font size
 viewparams\{2\} = []; \% The edges will be found by the algorithm
 viewparams\{3\} = 60; \% Those edges that appear at least 60 times
 viewparams\{4\} = 2;  \% Structures that have at least two edges
 viewparams\{5\} = 'correlation'; \% Similarity measure used to group edges is correlation
 [run\_structures,maxgen,nruns] = ReadStructures('ProteinStructsExR.txt',20);
 [results] = ViewStructures(run\_structures,20,maxgen,nruns,  'viewmatrix\_method',
 'ViewDendroStruct',viewparams);  \UnnumLine \end{programf}
 \label{ex:ViewDendroStruct} \end{example}

  Figure~\ref{fig:DendroVis} shows the  dendrogram of the most frequent edges (they appear
  at least $20$ times) in
   the structures  learned in several runs of EBNA-Exact
     and that satisfy the condition that the structures have at least
     two edges. These are the same edges shown in
     Figure~\ref{fig:PCVis} using the parallel coordinate
     representation.


\subsubsection{\texttt{ViewGlyphStruct} method}

 A Glyph is a visual representation of a piece of data where the attributes of a graphical
  entity (e.g.  shape, size, color, and position) are dictated by one or more attributes of a data record. The placement or layout of glyphs on a display can communicate
significant information regarding the data values themselves as well
as relationships between data points \cite{Ward:2002}. Among the
best known types of glyph graphs are star glyphs and Chernoff
glyphs \cite{Chernoff:1973}.  Glyphs
  have been extensively used as a visualization technique for many problems \cite{Lee_et_al:2003,Wittenbrink_et_al:1996}

Method \texttt{ViewGlyphStruct} displays  the glyph representation
of a subset of edges for a user defined set of runs and generations.
If none of the selected edges is included (the structure learned at
run $i$, generation $j$) a dot appears as the glyph representation.
This representation is useful to detect common substructures in the
models and to observe the complexity of the models learned along the
generations. Example~\ref{ex:ViewGlyphStruct} shows an example of
the application of the \texttt{ViewGlyphStruct} method.


\begin{example}[Application of \texttt{ViewGlyphStruct} method]
\begin{figure*}
  \begin{center}
   {{\includegraphics[width=9.0cm]{ViewGlyphEx.eps}}}
     \caption{Glyph visualization of the most frequent edges in the structures
     learned by EBNA-Exact.}
    \label{fig:GlyphVis}
    \end{center}
\end{figure*}

 \begin{programf}
 viewparams\{1\} = [14];
 viewparams\{2\} = results{3}; \% List of edges. From  results\{3\}, previous example.
 viewparams\{3\} = [1:10];     \% Selected runs  that will be  inspected
 viewparams\{4\} = [1:13];     \% Selected generations  that will be inspected
 [run\_structures,maxgen,nruns] = ReadStructures('ProteinStructsExR.txt',20);
 [newresults] = ViewStructures(run\_structures,20,maxgen,nruns, 'viewmatrix\_method',
 'ViewGlyphStruct',viewparams);  \UnnumLine \end{programf}
 \label{ex:ViewGlyphStruct} \end{example}

  Figure~\ref{fig:GlyphVis} shows the glyphs of the substructures formed by the
  most frequent edges found in  the previous two examples.
   However, in this case only the
  first $10$ runs and $13$ generations are considered to reduce the
  number of structures on display.


\section{Function approximation module}

   The goal of the function approximation module is to implement
   methods for using and validating the probabilistic models learned
   by EDAs for approximating functions. This research trend has
   received an increasing attention in the field of EDAs \cite{Brownlee_et_al:2008,Sastry_et_al:2004,Sastry_et_al:2006,Shakya:2006,Shakya_et_al:2005a}.
    In this
   section, we discuss the question of function approximation with models
   learned by EDAs, and then we present the methods implemented in MATEDA-2.0 for
   this purpose.

   \subsection{Probabilistic models of (possibly multi-objective) fitness functions}

   Probabilistic models of the fitness functions
   can be useful in different situations:

  \begin{itemize}
    \item To create surrogated functions that help to diminish the number of evaluations for costly functions \cite{Lima_et_al:2006,Orriols_et_al:2007,Sastry_et_al:2004,Sastry_et_al:2006,Shakya:2006,Shakya_et_al:2005a}.
    \item To obtain models of black box optimization problems for
    which an analytical expression of the fitness function is not
    available.
    \item To unveil and extract problem information that is hidden in the
    original formulation of the function or optimization problem \cite{Brownlee_et_al:2008}.
    \item To design improved (local) optimization procedures based on
    the model structure \cite{Butz_et_al:2006,Pereira_et_al:2000,Rivera_and_Santana:2000,Sastry_et_al:2006}.
  \end{itemize}

An important questions is: Given two probabilistic models of a
fitness function, which of the models is better? Notice that in this
case, models are not intended to be evaluated in terms of its
accuracy to represent the selected set (i.e. maximize the likelihood
of the data).  Instead, we would like to use them as sufficiently
general \emph{predictors} of the fitness function, or at least of
the fitness function of sufficiently good solutions. For
multi-objective problems, we can think of multi-models, each model
representing a different objective.

We make a number of assumptions. First, the models have been already
constructed. Actually, they are a byproduct of the EDA search.
Otherwise, other machine learning techniques devised for function
approximation could be used. The second assumption is that it is
possible to generate a sample of solutions for which the objectives
can be evaluated.  We then separate the problem of finding good
predictors in two subproblems:

\begin{enumerate}
  \item Which measures should be used to evaluate the models using the solutions (with their respective objective
  evaluations)?
  \item Which criteria should be taken into account to generate the solutions used to evaluate the models?
\end{enumerate}


Different criteria can be employed to compare the models:

\begin{itemize}
  \item Correlation between the probabilities assigned by the models to
  the solutions and their fitness values \cite{Brownlee_et_al:2008}.
  \item Sum of the probabilities assigned to the solutions \cite{Santana:2002}.
  \item Entropy of the model.
  \item Expected fitness value of the model, i.e. $\sum_x p(x)f(x)$.
\end{itemize}

These criteria will be closely related to the way in which the
solutions are selected. Let us suppose the number of solutions to be
generated is $k$. We identify as relevant the following procedures
to generate them:

\begin{itemize}
  \item The $k$ solutions are randomly generated \cite{Brownlee_et_al:2008}.
  \item The $k$ solutions correspond to the best known values (for a single objective) of the function
or are the selected solutions \cite{Santana:2002,Brownlee_et_al:2008}.
  \item Solutions correspond to the $k$ most probable configurations (MPC) of
  the model \cite{Echegoyen_et_al:2009}.
\end{itemize}

 For random solutions, we can compare different models in terms of
 the correlations or the expected fitness values. In \cite{Brownlee_et_al:2008}, the
 analysis of the correlation has been successfully
employed to analyze the fitness modeling capabilities of EDAs based on Markov networks.
  The entropy of the model can also be  used. A model that assigns the same or very
 similar probability to all the solutions is of scarce interest.

 If the best known solutions are used to evaluate the models, the previous criteria can be used.
 However, care must be taken to detect the phenomenon of overfitting since it may have a harmful
 effect in EDAs \cite{Lima_et_al:2008,Santana:2002,Santana:2005}.
 This can be done by computing the total probability assigned to the best solutions \cite{Santana:2002}.
 If the sum of the probabilities given by the
 model to the $k$ best solutions is very high (e.g. $\sum_x p(x) =
 0.9$) then we can assume its capacity of generalization is limited.

The $k$ most probable configurations can be computed using
algorithms that employ abductive inference and dynamic programming
as those presented in
\cite{Hoens_et_al:2007,Nilsson:1998,Yanover_and_Weiss:2004}. Most
probable configurations have been used in different contexts in EDAs
\cite{Echegoyen_et_al:2008a,Hoens_et_al:2007,Mendiburu_et_al:2007a,Soto:2003}.



\subsection{MATEDA-2.0 methods for fitness function approximation}




 The following methods are implemented in MATEDA-2.0:

\begin{itemize}
  \item  \texttt{BN\_Pop\_Prob}: Computes the probabilities given to a set of solutions by a
  Bayesian network.
  \item  \texttt{BN\_Fitness\_Corr}: Computes the correlation between the probabilities given to the solutions by a
   network and the fitness values of the  solutions.
  \item  \texttt{Find\_kMPEs}:  Given a Bayesian network, find the $k$ most probable
              configurations of the network.
  \item \texttt{Find\_Fitness\_Approx}:   Finds the probabilistic model with the
  highest correlation for each of the objectives in the given population (e.g. a Pareto set
  approximation).
\end{itemize}

The implementation of several of the methods discussed in the
previous section (e.g. generation of a random population,
computation of the sum of probabilities, computation of the entropy,
etc.) is straightforward in Matlab.

The \texttt{Find\_kMPEs} method is essentially the algorithm
introduced by Nilsson \cite{Nilsson:1998}. The algorithm computes
the junction tree of the BN and apply max-propagation and dynamic
programming for finding the $k$ MPCs.  In \cite{Nilsson:1998}, two
schedules for finding the subsequent maxima are proposed. This
implementation corresponds to the first proposed schedule.
Currently, it only works for binary variables.

 \texttt{Find\_Fitness\_Approx} can be directly applied to select the set
 of models that best approximate (in terms of the correlation values) a Pareto
 set approximation.


\section{Conclusions}


   It follows a summary of the main characteristics of the
   MATEDA-2.0 implementation.

\begin{itemize}
  \item It can be used for the optimization of single and multi-objective problems.
  \item Highly modular implementation in which each EDA component (either added by the user
  or already included in MATEDA-2.0) is implemented as an independent program.
  \item Available implementations include learning and sampling methods of undirected
  graphical models and Bayesian networks for problems with discrete and continuous variables.
  \item Knowledge extraction and visualization module for extracting
  and displaying information from the probabilistic models learned
  and the populations generated by the EDA.
  \item A variety of seeding, local optimization, repairing, selection and
  replacement methods.
  \item Statistical analysis of different measures of the EDA evolution.
  \item Extended library of functions and testbed problems.
\end{itemize}


\subsection{EDAs that can be implemented with MATEDA-2.0}

By combining the different implementations of the EDA components already included in MATEDA-2.0, variants
 of the following EDAs can be implemented:

\begin{itemize}
  \item Univariate marginal distribution algorithm (UMDA)
  \cite{Muhlenbein_and_Paas:1996}.
  \item Factorized distribution algorithm (FDA)
  \cite{Muhlenbein_et_al:1999}.
  \item EDAs based on trees and forests (Tree-EDA) \cite{Baluja_and_Davies:1997,Pelikan_and_Muhlenbein:1999,Santana_et_al:2001br}.
  \item EDAs based on Bayesian networks, similar to those presented
  in \cite{Etxeberria_and_Larranhaga:1999,Muehlenbein_and_Mahnig:2001a,Pelikan_et_al:1999}.
  \item Markov chain estimation of distribution algorithm (Mk-EDA)
  \cite{Santana_et_al:2004}.
  \item EDAs based on univariate and multivariate Gaussian distributions \cite{Bosman_and_Thierens:2000,Bosman_and_Thierens:2000a,Larranhaga_et_al:1999,Larranhaga_et_al:2000}.
  \item EDAs based on mixtures of continuous distributions \cite{Bosman_and_Thierens:2002}.
  \item Markov optimization algorithm (MOA)  \cite{Shakya_and_Santana:2008}.
  \item Affinity propagation EDA (Aff-EDA)  \cite{Santana_et_al:2008}.
\end{itemize}


In this paper we have presented MATEDA-2.0, a suite of programs in Matlab for optimization using EDAs.
We expect that these programs could help to find new applications of EDAs
to practical problems. The knowledge-extraction and visualizations methods introduced should
be useful in extending
 the uses of probabilistic modeling in optimization, particularly for revealing
 unknown information in black-box optimization problems.  In the future, we also intend
to incorporate new methods for the treatment of highly complicated, mixed, constrained, and other
 difficult problems \cite{Santana_et_al:2009b}.


\bibliographystyle{abbrv}
\bibliography{../NewThesis/Thesbib}

\end{document}


 The fitness
values of selected population are passed as parameters since some
learning algorithms

MATEDA-2.0 includes several procedures that
   allow the analysis and visualization of the information gathered during the evolution,
   particularly of the models structures and parameters.

To visualize the structures of the Bayesian models learned the
following commands from the Bayesian can be used.

\begin{programf}
 figure
 draw_graph(bnet.dag);
\end{programf}




Off-lattice models do not follow a given lattice topology. Instead,
the 2d or 3d coordinate in the real axis define the positions of the
residues. Among the off-lattice models with known lowest energy
states is the AB model  \cite{Stillinger_et_al:1993} where A stands
from hydrophobic and B from the polar residue. The distances between
consecutive residues along the chain are held fixed to $b=1$, while
non-consecutive residues interact through a modified Lennard-Jones
potential. There is an additional energy contribution from each
angle between successive bonds. We only consider Fibbonacci
sequences defined recursively.

Angles are represented as continuous values in $[0, 2\pi]$. We look
for the set of angles that defines a the optimal off-lattice
configuration. The univariate Gaussian EDA is implemented in folder
\texttt{GaussProtein}. The file \texttt{RunGaussianEDA.m} explains
the way solutions are codified. Other sequence configurations can be
added to this file.


@InProceedings{Orriols_et_al:2007,
  author = {Orioles-Puis and E. Bernard\'o-Manilla and Kumara Pastry and David E. Goldberg},
  title = "Substructures surrogates for learning decomposable classification problems: implementation and first results",
  booklice =    "Proceedings of the Genetic and Evolutionary  Computation Conference {GECKO}-2007 companion",
  pages = "2875--2882",
  year =         "2007",
  editor =       "D. {Thieves et al.}",
  address =      "London, UK",
  publisher = {ACE Press},
 }



@InProceedings{Lima_et_al:2006,
  author =       "C. F. Lima and M. Pelikan and K. Sastry and M. V. Butz and D. E. Goldberg and F. G. Lobo",
  title     = {Substructural neighborhoods for local search in the {B}ayesian optimization algorithm},
  editor    = {Thomas Philip Runarsson and  Hans-Georg Beyer and Edmund K. Burke and Juan J. Merelo Guerv{\'o}s and
               L. Darrell Whitley and Xin Yao},
  booktitle     = {Proceedings of the Parallel Problem Solving from Nature Conference (PPSN IX)},
  publisher = {Springer Verlag},
  series    = {Lecture Notes in Computer Science},
  volume    = {4193},
  year      = {2006},
  isbn      = {3-540-38990-3},
  pages     = {232-241},
}

@InProceedings{Vinko_et_al:2007,
  author =       "Tamas Vinko and Dario Izzo and Claudio Bombardelli",
  booktitle     = {Proceedings of 58th International Astronautical Congress},
  title     = {Benchmarking different global optimisation techniques for
preliminary space trajectory design},
  year      = {2007},
 }

@InProceedings{Posik:2008,
 author = {P. P\^os\'ik},
 title = {Preventing premature convergence in a simple {EDA} via global step size setting},
 booktitle = {Parallel Problem Solving from Nature - PPSN X},
 editor = {G. Rudolph and T. Jansen and S. Lucas and C. Poloni and N. Beume},
 year = {2008},
 series =  "Lecture Notes in Computer Science",
 publisher = {Springer Verlag},
 address = {Dortmund,Germany},
 volume = {5199},
 pages = {549--558},
}
